% \documentclass[a4paper,12pt]{article}{{{
\documentclass[a4paper,11pt]{article}
% -------------------------------------------------
% Set up the page margins.
% -------------------------------------------------
% \usepackage[text={6.5in,9.5in},centering]{geometry}
% \setlength{\topmargin}{-0.25in}
\usepackage[margin=1in]{geometry}

% -------------------------------------------------\left( \left[ \frac{a-d}{2}\right]q + p \right)
% Load Packages here
% -------------------------------------------------
%\usepackage{hyperref}
\usepackage{graphicx,url,subfig}
\RequirePackage[OT1]{fontenc}
\RequirePackage{amsthm,amsmath,amssymb,amscd}
\RequirePackage[numbers]{natbib}
\RequirePackage[colorlinks,citecolor=blue,urlcolor=blue]{hyperref}
\usepackage{graphics}
\usepackage{enumerate}
\usepackage{datetime}
\usepackage{etex}
% \usepackage[notcite,notref,color]{showkeys}
\usepackage[normalem]{ulem} % either use this (simple) or
\usepackage{soul} % use this (many fancier options)
\makeatother
\numberwithin{equation}{section}
\allowdisplaybreaks
\usepackage{mathrsfs}


% -------------------------------------------------
% check references
% -------------------------------------------------
% \usepackage{refcheck}
% \usepackage[notref,notcite]{showkeys}
% \usepackage[notcite]{showkeys}
% \usepackage{showkeys}

% -------------------------------------------------
% Environments here
% -------------------------------------------------
\theoremstyle{plain}
\newtheorem{theorem}{Theorem}[section]
\newtheorem{corollary}[theorem]{Corollary}
\newtheorem{lemma}[theorem]{Lemma}
\newtheorem{proposition}[theorem]{Proposition}
\newtheorem{exercise}{Exercise}
\theoremstyle{definition}
\newtheorem{definition}[theorem]{Definition}
\newtheorem{hypothesis}[theorem]{Hypothesis}
\newtheorem{assumption}[theorem]{Assumption}

\newtheorem{qu}{Question}

\newtheorem{remark}[theorem]{Remark}
\newtheorem{conjecture}[theorem]{Conjecture}
\newtheorem{example}[theorem]{Example}

% -------------------------------------------------
%  Define short hand symbols.
% -------------------------------------------------
\newcommand{\E}{\mathbb{E}}
\newcommand{\W}{\dot{W}}
\newcommand{\B}{\textbf{B}}
\newcommand{\D}{\mathbb{D}}
\newcommand{\ud}{\ensuremath{\mathrm{d} }}
\newcommand{\Ceil}[1]{\left\lceil #1 \right\rceil}
\newcommand{\Floor}[1]{\left\lfloor #1 \right\rfloor}
\newcommand{\sgn}{\text{sgn}}
\newcommand{\Lad}{\text{L}_{\text{ad}}^2}
\newcommand{\SI}[1]{\mathcal{I}\left[#1 \right]}
\newcommand{\SIB}[2]{\mathcal{I}_{#2}\left[#1 \right]}
\newcommand{\Indt}[1]{1_{\left\{#1 \right\} }}
\newcommand{\LadInPrd}[1]{\left\langle #1 \right\rangle_{\text{L}_\text{ad}^2}}
\newcommand{\LadNorm}[1]{\left|\left|  #1 \right|\right|_{\text{L}_\text{ad}^2}}
\newcommand{\Norm}[1]{\left\|  #1   \right\|}
\newcommand{\Ito}{It\^{o} }
\newcommand{\Itos}{It\^{o}'s }
\newcommand{\spt}[1]{\text{supp}\left(#1\right)}
\newcommand{\InPrd}[1]{\left\langle #1 \right\rangle}
\newcommand{\mr}{\textbf{r}}
\newcommand{\Ei}{\text{Ei}}
\newcommand{\arctanh}{\operatorname{arctanh}}
\newcommand{\ind}[1]{\mathbb{I}_{\left\{ {#1} \right\} }}

\DeclareMathOperator{\esssup}{\ensuremath{ess\,sup}}

\newcommand{\steps}[1]{\vskip 0.3cm \textbf{#1}}
\newcommand{\calB}{\mathcal{B}}
\newcommand{\calD}{\mathcal{D}}
\newcommand{\calE}{\mathcal{E}}
\newcommand{\calF}{\mathcal{F}}
\newcommand{\calG}{\mathcal{G}}
\newcommand{\calK}{\mathcal{K}}
\newcommand{\calH}{\mathcal{H}}
\newcommand{\calI}{\mathcal{I}}
\newcommand{\calL}{\mathcal{L}}
\newcommand{\calM}{\mathcal{M}}
\newcommand{\calN}{\mathcal{N}}
\newcommand{\calO}{\mathcal{O}}
\newcommand{\calT}{\mathcal{T}}
\newcommand{\calP}{\mathcal{P}}
\newcommand{\calR}{\mathcal{R}}
\newcommand{\calS}{\mathcal{S}}
\newcommand{\bbC}{\mathbb{C}}
\newcommand{\bbN}{\mathbb{N}}
\newcommand{\bbP}{\mathbb{P}}
\newcommand{\bbZ}{\mathbb{Z}}
\newcommand{\myVec}[1]{\overrightarrow{#1}}
\newcommand{\sincos}{\begin{array}{c} \cos \\ \sin \end{array}\!\!}
\newcommand{\CvBc}[1]{\left\{\:#1\:\right\}}
\newcommand*{\one}{ {{\rm 1\mkern-1.5mu}\!{\rm I} }}
\newcommand{\RR}{\mathbb{R}}
\newcommand{\R}{\mathbb{R}}
\newcommand{\myEnd}{\hfill$\square$}
\newcommand{\ds}{\displaystyle}
\newcommand{\Shi}{\text{Shi}}
\newcommand{\Chi}{\text{Chi}}
\newcommand{\Daw}{\ensuremath{\mathrm{daw} }}
\newcommand{\Erf}{\ensuremath{\mathrm{erf} }}
\newcommand{\Erfi}{\ensuremath{\mathrm{erfi} }}
\newcommand{\Erfc}{\ensuremath{\mathrm{erfc} }}
\newcommand{\He}{\ensuremath{\mathrm{He} }}

\def\crtL{\rho}

\newcommand{\conRef}[1]{(#1)}

\DeclareMathOperator{\Lip}{\mathit{L}}
\DeclareMathOperator{\LIP}{Lip}
\DeclareMathOperator{\lip}{\mathit{l}}

\DeclareMathOperator{\Vip}{\overline{\varsigma}}
\DeclareMathOperator{\vip}{\underline{\varsigma}}
\DeclareMathOperator{\vv}{\varsigma}

% % Set up mdframe for questions.
\newcounter{NQ}
\newcounter{LQ}
\newcounter{ZQ}
\usepackage{xcolor}
\usepackage[framemethod=tikz]{mdframed}
\mdfsetup{%
	middlelinecolor=lightgray,
	middlelinewidth=2pt,
	backgroundcolor=red!20,
	roundcorner=10pt}
\newenvironment{NickQuestion}[1][] %
{\refstepcounter{NQ}
	\begin{mdframed}[backgroundcolor=lightgray]\textbf{Nick's Question \theNQ \: #1~~}~~~\noindent}%
	{\end{mdframed}}
\newenvironment{LeQuestion}[1][] %
{\refstepcounter{LQ}
	\begin{mdframed}[backgroundcolor=red!10]\textbf{Le's Question \theLQ \: #1~~}~~~\noindent}%
	{\end{mdframed}}
\newenvironment{ZhijianQuestion}[1][] %
{\refstepcounter{ZQ}
	\begin{mdframed}[backgroundcolor=blue!10]\textbf{Zhijian's Question \theZQ \: #1~~}~~~\noindent}%
	{\end{mdframed}}
\numberwithin{NQ}{section}
\numberwithin{LQ}{section}
\numberwithin{ZQ}{section}

\usepackage[notref,notcite]{showkeys}

\title{Nick Eisenberg's Meeting Notes}
\author{Le Chen, Nick Eisenberg, Zhijian Wu}
\date{October 2019}

\usepackage{natbib}
\usepackage{graphicx}
% }}}
\begin{document}
	\section{Introduction}
	We study the interpolated stochastic heat and wave equation (ISHWE) for spatial dimension on $\R^d$,
	\begin{equation}
	\label{E:SPDE}
	\begin{cases}
	\left(\partial^\beta_t + \frac{\nu}{2}(-\Delta)^{\alpha/2}\right)\: u(t,x) = I^\gamma_t \left[\sigma( u(t,x)) \: \dot W(t,x) \right] & \text{$x\in \R^d$, $t>0$} \\
	u(0,\cdot) = u_0                                                                                                            & \beta \in (0,1]               \\
	u(0,\cdot) = u_0, \quad \partial_t u(0,\cdot) = v_0                                                                            & \beta \in (1,2).
	\end{cases}
	\end{equation}
	
	
	\section{Preliminaries}
For any $t \in [0,T]$ and $x \in \R^d$, we write (assuming $\beta \in (1,2))$
\[
u(t,x) = J_0(t,x) + I(t,x)
\]
	where
	\begin{align*}
	J_0(t,x) &=\begin{cases} 
	[Z(t,\cdot)*u_0](x) & \beta \in (0,1]
	\\	[Z(t,\cdot)*v_0](x) + \frac{\partial}{\partial t}[Z(t,\cdot)*u_0](x) & \beta \in (1,2)
	\end{cases}
	\\	I(t,x) &= \int_0^t \int_{\R^d} Y(t-s,x-y)\sigma(u(s,y)) W(\ud s, \ud y).
	\end{align*}

The function $Y$ is defined through the Fox-H function and its Fourier transform can be shown to be given as 
	\[
		\mathcal{F}Y(t,\cdot)(\xi) = t^{\beta + \gamma-1}E_{\beta. \beta +\gamma}(-2^{-1}\nu t^\beta |\xi|^\alpha)
	\]
	
	\textcolor{blue}{The following is for space-time white noise. For colored noise we need $\rho(\lambda) >0$ and $2\alpha > d$.}
	
	\section{Space-time White Noise}
	In this section we study \eqref{E:SPDE} for driven by space-time white noise. To establish the existence and uniqueness of random field solutions in the to \eqref{E:SPDE}, we will need to assume Dalang's condition:
	\begin{equation}
	\label{C:Dalang}
	\int_0^t \ud s \int_{\R^d} \ud y |Y(s,y)|^2 < \infty, \quad \text{for all } t>0,
	\end{equation} 
	which is equivalent to
	\begin{equation}
	\label{E:DalangEquiv}
	d < 2\alpha + \frac{\alpha}{\beta}\min\{2\gamma -1, 0\} =: \Theta,
	\end{equation}
	which is further equivalent to 
	\begin{equation}
	\label{E:DalangEquiv2}
	\rho(d) > 0 \quad \text{and} \quad d < 2\alpha,
	\end{equation}
	where 
	\begin{equation}
	\label{D:theta}
	\rho(x) := 2\beta + 2\gamma -1 -\beta x / \alpha.
	\end{equation} 
	\subsection{Globally Lipschitz Coefficient}
	In this section, we assume that the coefficient $\sigma$ is globally Lipschitz and therefore satisfies the following inequality:
	\[
	|\sigma(x)| \le c(\sigma) + L(\sigma)|x|,
	\] 
	where $c(\sigma) = \sigma(0)$ and $L(\sigma)$ is the Lipschitz constant of $\sigma$. We assume that $L(\sigma)>0$. We now state an existence and uniqueness result proven in \cite{CHN19}.
	
	\begin{theorem}
		\label{T:SolveabilityGenLip}
		Under \eqref{C:Dalang}, the SPDE \eqref{E:SPDE} has a unique (in the sense of versions) random field solution $\{u(t,x) : (t,x) \in (0,\infty) \times \R^d \}$ if the initial data are such that 
		\begin{equation}
		\hat{C}_t := \sup_{t\in[0,T],\; x\in \R^d} |J_0(t,x)| < \infty.
		\end{equation}	
		\textcolor{blue}{Moreover, the following statements are true:
			\begin{enumerate}
				\item $(t,x) \mapsto u(t,x)$ is $L^P(\Omega)$-continuous for all $p \ge 2$,
				\item For all even integers $p \ge 2$, all $t \ge 0$ and $x,y\in\R^d$,
				\begin{equation}
				\Norm{u(t,x)}_p^2 \le 2J_0^2(t,x) + \left[ \bar{\varsigma}^2 + 2 \hat{C}_t^2 \right]\exp(C \lambda^{2/(1-\sigma)}t).
				\end{equation}
		\end{enumerate}}
		
	\end{theorem}
	\begin{lemma}{\normalfont \cite[Theorem 4.1 and Lemma 5.5]{CHN19}}
		\label{L:1and2normYZ}
		Assume that $d < 2\alpha$, $\beta \in (0,2)$, and $\gamma \ge 0$. Then 
		\[
		\int_{\R^d} Y(t,x) \ud x = t^{\beta+\gamma -1 }
		\]
		and
		\[
		\int_{\R^d} Y^2(t,x) \ud x = C_{\#} t^{2(\beta+\gamma-1)-d\beta/\alpha} = C_{\#} t^{\rho -1} ,
		\]
		for all $t > 0$, where $\rho = \rho(d)$ and 
		\[
		C_{\#} := \frac{2}{\Gamma(d/2)(4 \pi )^{d/2}} \int_0^\infty u^{d-1}E^2_{\beta,\beta + \gamma}(-u^\alpha) \ud u.
		\]
		Now further assume that $\beta \in (1,2)$, then 
		\[
		\int_{\R} Z(t,x) \ud x= t^{\lceil \beta \rceil -1} = t. 
		\]
	\end{lemma}
	\begin{remark}
		When $\alpha = \beta = 2$,  $\gamma = 0$ and $d=1$ then Lemma \ref{L:1and2normYZ} gives us that $\int_{\R} G^2(t,x) \ud x = t/2$ which coincides with the $L^2(\R)$ norm of the wave kernel. 
	\end{remark}
	In the proof of the next proposition, we will use the following two facts:
	\begin{equation}
	\label{E:SUPidentities}
	\sup_{t\ge0} t^ke^{-at} = k^k(ea)^{-k} \text{ for }k \in \mathbb{N} \text{ and } \sup_{t \ge 0} \int_0^t s e^{-as} \ud s = a^{-2} \text{ for } a>0.
	\end{equation}
	
	\begin{proposition}{\normalfont\cite[Proposition 3.2]{SoleMillet2021}}
		\label{P:N(u)Bdd}
		Let $u_0$ and $v_0$ be Borel functions satisfying $\Norm{u_0}_\infty + \Norm{v_0}_\infty < \infty$. Suppose that $d < 2 \alpha$ and $\rho = \rho(d) > 0$, where $\rho(x)$ is defined in \eqref{D:theta}. Then there exists a universal constant $K:=K(\alpha, \beta, \gamma,d) >0$ such that for any 
		\[
		a > \left( 8 L^2(\sigma)K^2\right)^{1/\rho}.
		\]
		and 
		\[
		p \in \left[ 2, \frac{a^{\rho}}{4L^2(\sigma)K^2} \right].
		\]
		we have that 
		\begin{equation}
		\label{E:Nbdd}
		N_{a,p}(u)  \le 2 \mathcal{T}_0 +  \frac{c(\sigma)}{L(\sigma)}
		\end{equation}	
		where
		\begin{align*}
		\mathcal{T}_0 = \begin{cases} 
		(\lceil \beta \rceil - 1)^{\lceil \beta \rceil -1}(ea)^{1-\lceil \beta \rceil} \Norm{u_0}_\infty  & \beta \in (0,1]
		\\ (\lceil \beta \rceil - 1)^{\lceil \beta \rceil -1}(ea)^{1-\lceil \beta \rceil} \Norm{v_0}_\infty + \Norm{u_0}_\infty & \beta \in (1,2).
		\end{cases} 
		\end{align*}
		Moreover for any $T>0$,
		\begin{equation}
		\label{E:Pnormtbdd}
		\sup_{x\in\R^d} \mathbb{E}(|u(t,x)|^p) \le \exp\left( a p t  \right) \left[ 2 \mathcal{T}_0 + \frac{ c(\sigma) }{L(\sigma)} \right]^p, \; t\in[0,T].
		\end{equation}
	\end{proposition}
	\begin{proof}
		First we fix an $a>0$ and $p \in [2,\infty)$. Using Lemma \ref{L:1and2normYZ} we see that,
		\begin{align*}
		|J(t,x)|  &= \begin{cases} 
		\left| (Z(t,\cdot) * u_0)(x)  \right| & \beta \in (0,1]
		\\	\left| (Z(t,\cdot) * v_0)(x) + \frac{\partial}{\partial t}(Z(t,\cdot) * u_0)(x) \right| & \beta \in (1,2)
		\end{cases}
		\\& \le \begin{cases}
		t^{\lceil \beta \rceil -1}\Norm{u_0}_\infty & \beta \in (0,1]
		\\ t^{\lceil \beta \rceil -1}\Norm{v_0}_\infty + \Norm{u_0}_\infty & \beta \in (1,2)
		\end{cases} 
		\end{align*}
		Now using \eqref{E:SUPidentities}, we see that 
		\begin{align}
		N_{a,p}(J) &\le \begin{cases} 
		\sup_{t\ge 0}\sup_{x\in\R}  t^{\lceil \beta \rceil -1}\Norm{u_0}_\infty e^{-at} & \beta \in (0,1]
		\\ \sup_{t\ge 0}\sup_{x\in\R} \left(  t^{\lceil \beta \rceil -1}\Norm{v_0}_\infty + \Norm{u_0}_\infty \right) e^{-at} & \beta \in (1,2)
		\end{cases} \nonumber
		\\ &= \begin{cases} 
		(\lceil \beta \rceil - 1)^{\lceil \beta \rceil -1}(ea)^{1-\lceil \beta \rceil} \Norm{u_0}_\infty  & \beta \in (0,1]
		\\ (\lceil \beta \rceil - 1)^{\lceil \beta \rceil -1}(ea)^{1-\lceil \beta \rceil} \Norm{v_0}_\infty + \Norm{u_0}_\infty & \beta \in (1,2). 
		\end{cases} \label{E:Jbdd}
		\end{align}
		
		Now we find an upper bound for $N_{a,p}(I)$. By simultaneously applying the Burkholder-Davies-Gundy inequality and triangle inequality, we get that  
		\begin{align*}
		\Norm{I(t,x)}_p^2 &\le 4p \Norm{\int_0^t \ud s \int_{\R^d}\ud y \; Y^2(t-s,x-y)\sigma^2(u(s,y))  }_{p/2}
		\\&\le 4p \int_0^t \ud s \int_{\R^d} \ud y \; Y^2(t-s,x-y)\Norm{\sigma^2(u(s,y))}_{p/2}.
		\end{align*}
		Now by applying Lemma \ref{L:1and2normYZ} and the Lipschitz property of $\sigma$ we get that 
		\begin{align*}
		\Norm{I(t,x)}_p^2 &\le 8pc^2(\sigma)C_{\#}\int_0^t (t-s)^{\rho -1} \ud s + 8pL^2(\sigma)N^2_{a,p}(u)\int_0^t \ud s \int_{\R^d} \ud y \; e^{2as} Y^2(t-s,x-y).
		\end{align*} 
		Now under Dalang's condition \eqref{E:DalangEquiv2}, the first integral above can be calculated as 
		\[
		\int_0^t (t-s)^{\rho -1} \ud s = \frac{ t^{\rho} }{\rho}.
		\]
		By applying Lemma \ref{L:1and2normYZ}, we can calculate the second integral as 
		\[
		\int_0^t \ud s \int_{\R} \ud y \; e^{2as} Y^2(t-s,x-y) = C_{\#}\int_0^t \ud s \; e^{2as}(t-s)^{\rho - 1}.
		\]
		Putting this together gives us that 
		\begin{align*}
		\Norm{I(t,x)}_p^2 &\le 8pc^2(\sigma)C_{\#} \frac{ t^{\rho} }{\rho} + 8pL^2(\sigma)N^2_{a,p}(u) C_{\#}\int_0^t \ud s \; e^{2as}(t-s)^{\rho - 1}.
		\end{align*}
		Now by multiplying both sides by $e^{-2at}$, we see that
		\begin{align*}
		\Norm{I(t,x)}_p^2e^{-2at} &\le \frac{8pc^2(\sigma)C_{\#} }{\rho} \; t^{\rho} e^{-2at} 
		\\&\qquad\qquad + 8pL^2(\sigma)N^2_{a,p}(u) C_{\#} \int_0^t \ud s \; e^{-2a(t-s)}(t-s)^{\rho - 1}.
		\end{align*}
		We now do a change of variables and again recall the definition of the Gamma function to see that 
		\begin{align*}
		\Norm{I(t,x)}_p^2e^{-2at}  &\le \frac{8pc^2(\sigma)C_{\#} }{\rho} \; t^{\rho} e^{-2at} + 8pL^2(\sigma)N^2_{a,p}(u) C_{\#} \int_0^\infty \ud u \; e^{-2au}(u)^{\rho -1} \\
		& = \frac{8pc^2(\sigma)C_{\#} }{\rho} \; t^{\rho} e^{-2at} + 8pL^2(\sigma)N^2_{a,p}(u) \frac{C_{\#}\Gamma\left(\rho \right)}{ (2a)^{\rho}}
		\end{align*}
		where the last line is true if Dalang's condition \cite{CHN19} holds. If we now apply \eqref{E:SUPidentities} and also recall that $\sqrt{a+b} \le \sqrt{a} + \sqrt{b}$, then we see that 
		\begin{align}
		N_{a,p}(I) &\le  c(\sigma)\sqrt{\frac{8pC_{\#}  }{\rho}} \left(\frac{\rho}{2ea}\right)^{\rho/2} +  L(\sigma)N_{a,p}(u) \sqrt{\frac{8pC_{\#} \Gamma\left( \rho \right)}{(2a)^\rho} }. \label{E:SoleMil3.9}
		\end{align}
		To simplify things, we now let 
		\begin{align*}
		K&:= \max\Bigg\{\sqrt{\frac{8C_{\#}  }{\rho}} \left(\frac{\rho}{2e}\right)^{\rho/2}, \sqrt{\frac{8C_{\#} \Gamma\left( \rho \right) }{2^{\rho}}}  \Bigg\}
		\end{align*}
		then,
		\begin{align}
		\label{E:Ibdd}
		N_{a,p}(I) &\le\frac{K \sqrt{p}}{a^{\rho/2}}\left( c(\sigma) + L(\sigma)N_{a,p}(u) \right).
		\end{align}
		Now by combining \eqref{E:Jbdd} and \eqref{E:Ibdd}, we see that 
		\begin{align}
		N_{a,p}(u) & \le \begin{cases} 
		(\lceil \beta \rceil - 1)^{\lceil \beta \rceil -1}(ea)^{1-\lceil \beta \rceil} \Norm{u_0}_\infty  & \beta \in (0,1]
		\\ (\lceil \beta \rceil - 1)^{\lceil \beta \rceil -1}(ea)^{1-\lceil \beta \rceil} \Norm{v_0}_\infty + \Norm{u_0}_\infty & \beta \in (1,2). 
		\end{cases} 
		\\& \quad \quad  + K\frac{\sqrt{p}c(\sigma)}{a^{\rho/2} }  + K  \frac{\sqrt{p}L(\sigma)}{a^{\rho/2}} N_{a,p}(u). \label{E:Nbddgen}
		\end{align}
		It is important to note that if we now restrict ourselves to the case of $\alpha = \beta = 2$, $d=1$ and $\gamma =0$, then \eqref{E:SoleMil3.9} and \eqref{E:Nbddgen} recover \cite[(3.9) and (3.10)]{SoleMillet2021} respectively. Now choose $a>0$ such that 
		\[
		a > \left( 8 L^2(\sigma)K^2\right)^{1/\rho}.
		\]
		This implies that the following interval is nonempty:
		\[
		\left[ 2, \frac{a^{\rho}}{4L^2(\sigma)K^2} \right].
		\]
		In addition, for any $p$ in this interval, we have that 
		\[
		\frac{\sqrt{p}L(\sigma)K}{a^{\rho/2}} \le \frac{1}{2},
		\]
		and we also have 
		\[
		\frac{\sqrt{p}}{a^{\rho/2}}  \le \frac{1}{2KL(\sigma)}.
		\]
		This along with \eqref{E:Nbddgen} implies that 
		\begin{align*}
		N_{a,p}(u) & \le 2 \begin{cases} 
		(\lceil \beta \rceil - 1)^{\lceil \beta \rceil -1}(ea)^{1-\lceil \beta \rceil} \Norm{u_0}_\infty  & \beta \in (0,1]
		\\ (\lceil \beta \rceil - 1)^{\lceil \beta \rceil -1}(ea)^{1-\lceil \beta \rceil} \Norm{v_0}_\infty + \Norm{u_0}_\infty & \beta \in (1,2)
		\end{cases} 
		\\& \quad \quad + \frac{c(\sigma)}{L(\sigma)}
		\end{align*}
		and this is precisely \eqref{E:Nbdd}. Note that \eqref{E:Pnormtbdd} follows immediately from \eqref{E:Nbdd}.  
	\end{proof}
	
	We will need to recall the following result proved in \cite{CHN19} in order to prove Proposition \ref{P:NormInc} below. 
	
	\begin{proposition}{\normalfont\cite[Proposition 2.1 and 2.2]{ChenHu21}}
		\label{P:YInc}
		Suppose that $\alpha \in (0,2]$, $\beta \in (0,2)$, $\gamma>0$, and \eqref{E:DalangEquiv} holds. Consider $\rho = \rho(d)$ which is defined above in \eqref{D:theta}. Then $Y(t,x) = Y_{\alpha, \beta, \gamma, d}(t,x)$ satisfies the following:
		\begin{enumerate}
			\item \label{P:YIncCase1}  Suppose that $0< \theta < (\Theta -d) \wedge 2$ where $\Theta$ is given in \eqref{E:DalangEquiv}, $0< t,r \le T$ for some $T>0$, $\beta \le 1$ and $\gamma \le \lceil \beta \rceil - \beta$. Then there exists some $C := C(\alpha, \beta, \gamma, \nu, d, \theta, T)$ such that
			\[
			\int_0^T \ud s \int_{\R^d} \ud y |Y(t-s,x-y) - Y(r-s,z-y)|^2 \le C \left(|t-r|^\rho + |x-z|^{\theta}\right).
			\]
			\item \label{P:YIncCase2} If we only assume that $\alpha \in (0,2]$, $\beta \in (0,2)$, $\gamma>0$, $0< t,r \le T$ for some $T>0$ and \eqref{E:DalangEquiv} holds then for $0< \theta < (\Theta -d) \wedge 2$ there exists some $C := C(\alpha, \beta, \gamma, \nu, d, \theta, T)$ such that
			\[
			\int_0^T \ud s \int_{\R^d} \ud y |Y(t-s,x-y) - Y(r-s,z-y)|^2 \le C \left(|t-r|^q + |x-z|^{\theta}\right),
			\]
			where 
			\[
			0 < q < \rho.
			\]
		\end{enumerate}
	\end{proposition}
	
	\begin{remark}
		In Proposition \ref{P:YInc} above, we may often ignore the dependence of $C$ on $\alpha, \beta, \gamma, \nu, d$ and just write $C(\theta, T)$. 
	\end{remark}
	
	\begin{proposition}
		\label{P:NormInc}
		Suppose that $u_0=1$ and that $v_0 =0$ so that $u(t,x) = 1 + I(t,x)$. Let $a>0$ and suppose that $0< \theta < (\Theta -d) \wedge 2$ and $\rho = \rho(d) > 0$. Then for any $x,z \in \R$, $p\ge 2$ and for any $T>0$ with $t\in [0,T]$,
		\begin{equation}
		\label{E:NormInc_N(u)}
		\frac{\Norm{u(t,x)-u(r,z)}_p}{\left(|t-r|^{q} + |x-z|^{\theta}\right)^{1/2}} \le 	C(p,\theta, T)\left[ \mathcal{M}_1 + \mathcal{M}_2 e^{aT} N_{a,p}(u)\right] ,
		\end{equation}
		where 
		\[
		\mathcal{M}_1 = \sqrt{p}c(\sigma) \quad \text{and} \quad\mathcal{M}_2 =\sqrt{p}L(\sigma),
		\]
		and
		\begin{enumerate}
			\item $0< q \le \rho$ under Case \ref{P:YIncCase1} of Proposition \ref{P:YInc},
			\item $0 < q < \rho$ under Case \ref{P:YIncCase2} of Proposition \ref{P:YInc}.
		\end{enumerate}
		Moreover, for any 
		\[
		a > \left( 8 L^2(\sigma)K^2\right)^{1/
			\rho}.
		\]
		and 
		\[
		p \in \left[ 2, \frac{a^{\rho}}{4L^2(\sigma)K^2} \right]
		\]
		we have that 
		\begin{equation}
		\label{E:NormInc_cL}
		\frac{\Norm{u(t,x)-u(r,z)}_p}{\left(|t-r|^{q} + |x-z|^{\theta}\right)^{1/2}} \le C(p,\theta, T)\left[ \mathcal{M}_1 + \mathcal{M}_2 e^{aT} \left( 2 +  \frac{c(\sigma)}{L(\sigma)} \right)\right].
		\end{equation}
	\end{proposition}
	
	\begin{proof}
		Note that because of the initial conditions, we see that $\Norm{u(t,x) - u(r,z)}_p = \Norm{I(t,x) - I(r,z)}_p$. We now apply Burkholder-Davies-Gundy inequality along with the triangle inequality, as was done in the proof of Proposition \ref{P:N(u)Bdd},  to see that
		\begin{align*}
		\Norm{I(t,x) - I(r,z)}_p^2 &\le 4p \Norm{ \int_0^T \ud s \int_{\R^d} \ud y \left| Y(t-s,x-y) - Y(r-s,z-y) \right|^2\sigma^2(u(s,y)) }_{p/2} \\
		&\le 4p \int_0^T \ud s \int_{\R^d} \ud y \left| Y(t-s,x-y) - Y(r-s,z-y) \right|^2  \Norm{\sigma^2(u(s,y)) }_{p/2}.
		\end{align*} 
		Now by applying the Lipchitz condition on $\sigma$ gives us that 
		\begin{align*}
		\Norm{I(t,x) - I(r,z)}_p^2 &\le 8p \int_0^T \ud s \int_{\R^d} \ud y \left| Y(t-s,x-y) - Y(r-s,z-y) \right|^2 \left( c^2(\sigma) + L^2(\sigma)\Norm{u(s,y)}_{p}^2\right) \\
		&\le 8p\Bigg[ c^2(\sigma)   \int_0^T \ud s \int_{\R^d} \ud y \left| Y(t-s,x-y) - Y(r-s,z-y) \right|^2 \\
		& \qquad+ L^2(\sigma) \int_0^T \ud s \int_{\R^d} \ud y \left| Y(t-s,x-y) - Y(r-s,z-y) \right|^2 \exp\left(2as\right)N_{a,p}^2(u)  \Bigg] \\
		&\le 8p\Bigg[ c^2(\sigma)   \int_0^T \ud s \int_{\R^d} \ud y \left| Y(t-s,x-y) - Y(r-s,z-y) \right|^2 \\
		& \qquad+ L^2(\sigma) \exp\left(2aT \right)N_{a,p}^2(u)  \int_0^T \ud s \int_{\R^d} \ud y \left| Y(t-s,x-y) - Y(r-s,z-y) \right|^2  \Bigg]. 
		\end{align*}
		We now apply Proposition \ref{P:YInc} to see that 
		\begin{align*}
		\Norm{I(t,x) - I(r,z)}_p^2 &\le 8p \bigg[ c^2(\sigma)C(\theta, T)\left(|t-r|^{q} + |x-z|^{\theta}\right)\\ 
		& \qquad +  L^2(\sigma) \exp\left(2aT \right)N_{a,p}^2(u) C(\theta, T)\left(|t-r|^{q} + |x-z|^{\theta}\right) \bigg],
		\end{align*}
		where $C(\theta, T)$ is as in Proposition \ref{P:YInc}. Now by taking the square root of both sides and applying the identity $\sqrt{x+y} \le \sqrt{x} + \sqrt{y}$ yields
		\begin{align*}
		\Norm{I(t,x) - I(r,z)}_p &\le 2\sqrt{2C(\theta, T)p} \Big[ c(\sigma) + L(\sigma) \exp\left(aT \right)N_{a,p}(u) \Big]\left(|t-r|^{q} + |x-z|^{\theta}\right)^{1/2}. 
		\end{align*}
		This proves \eqref{E:NormInc_N(u)}. It is clear that \eqref{E:NormInc_cL} follows directly from \eqref{E:NormInc_N(u)} and Proposition \ref{P:N(u)Bdd}. 
	\end{proof}
	
	\textcolor{blue}{The holder continuity in the next proposition needs to be checked. Also I need to change the $[-R,R]$ to $|x| \le R$ in the following proof so we can consider higher dimension than $d=1$.}
	
	\textcolor{blue}{Check the following Holder continuity with \cite{ChenHu21}.}
	
	\begin{proposition}
		\label{P:Holder}
		Suppose that $u_0$ and $v_0$ are as on Proposition \ref{P:NormInc} above. Then $u$ has a version, still denoted by $u$ that is $\eta_1$-H\"older continuous in time and $\eta_2$-H\"older continuous in space with $0 < \eta_1 < q/2$ and $0 < \eta_2 < \theta/2$ where $0 < \theta < (\Theta -d) \wedge 2$ and 
		\begin{enumerate}
			\item $0< q \le \rho$ under Case \ref{P:YIncCase1} of Proposition \ref{P:YInc},
			\item $0 < q < \rho$ under Case \ref{P:YIncCase2} of Proposition \ref{P:YInc}.
		\end{enumerate}
		Moreover, if we consider
		\[
		a > \left( 8 L^2(\sigma)K^2\right)^{1/\rho} \quad \text{and} \quad 	p \in \left[ 2, \frac{a^{\rho}}{4L^2(\sigma)K^2} \right],
		\]
		where $K$ is the constant from Proposition \ref{P:N(u)Bdd}, then we have that 
		\[
		\mathbb{E}\left[ \sup_{t\in[0,T],\; |x| \le R} |u(t,x)|^p \right] \le 2^{p-1} + C(p, \theta, T,R) \left[ \mathcal{M}_1^p + \mathcal{M}_2^p e^{apT} \left(2 +  \frac{c(\sigma)}{L(\sigma)}\right)^p\right].
		\]
	\end{proposition}
	
	\begin{proof}
		Proposition \ref{P:NormInc} implies that 
		\begin{equation}
		\Norm{u(t,x) - u(r,z)}_p \le C(p,\theta, T)\left[ \mathcal{M}_1 + \mathcal{M}_2e^{aT}N_{a,p}(u) \right]\left(|t-r|^{q} + |x-z|^{\theta}\right)^{1/2}.
		\end{equation}
		By Kolmogorov's continuity theorem, there exists a version of $u$, which we denote by $u$, that is $\eta_1$-H\"older continuous in time and $\eta_2$-H\"older continuous in space with with $0<\eta_1 < q/2$ and $0<\eta_2 < \theta/2$.
		
		Next, note that by triangle inequality,
		\begin{align*}
		|u(t,x)|^p &\le 2^{p-1}\left( \Norm{u_0}_{\infty}^p + |u(t,x)-u_0(x)|^p \right) \\
		&\le 2^{p-1}\left( \Norm{u_0}_{\infty}^p + \sup_{t\in[0,T],\; x\in[-R,R]}\left(|u(t,x)-u_0(x)|^p\right) \right) \\
		&= 2^{p-1}\left( 1 + \sup_{t\in[0,T],\; x\in[-R,R]}\left(|u(t,x)-1|^p\right) \right)
		\end{align*}
		Due to the continuity of $u$ on the compact set $[0,T] \times [-R,R]$, there exists a point $(t_0,x_0) \in [0,T] \times [-R,R]$ such that 
		\[
		\sup_{t\in[0,T],\; x\in[-R,R]}\left(|u(t,x)-1|^p\right) = |u(t_0,x_0)-1|^p.
		\]
		Thus, for any $(t,x) \in [0,T] \times [-R,R]$,
		\[
		|u(t,x)|^p \le 2^{p-1}\left( 1  + |u(t_0,x_0)-1|^p \right)
		\] 
		which in return implies that 
		\[
		\sup_{t\in[0,T],\; x\in[-R,R]} |u(t,x)|^p \le 2^{p-1}\left(1 + |u(t_0,x_0)-1|^p \right).
		\]
		Now by taking the expectation of both sides and applying Proposition \ref{P:NormInc} gives us 
		\begin{align*}
		\mathbb{E}\left[ \sup_{t\in[0,T],\; x\in[-R,R]} |u(t,x)|^p \right] &\le  2^{p-1} + 2^{p-1}C(p, \theta, T)^p\left((2T)^{q/2} + (2R)^{\theta /2}\right)^p\left[ \mathcal{M}_1 + \mathcal{M}_2 e^{aT} N_{a,p}(u)\right]^p \\
		&= 2^{p-1} + C(p, \theta, T,R) \left[ \mathcal{M}_1^p + \mathcal{M}_2^p e^{apT} N_{a,p}(u)^p\right].
		\end{align*}
		Now if we consider
		\[
		a > \left( 8 L^2(\sigma)K^2\right)^{1/\rho} \quad \text{and} \quad 	p \in \left[ 2, \frac{a^{\rho}}{4L^2(\sigma)K^2} \right],
		\]
		then by Proposition \ref{P:N(u)Bdd}, we have that 
		\[
		\mathbb{E}\left[ \sup_{t\in[0,T],\; x\in[-R,R]} |u(t,x)|^p \right] \le 2^{p-1} + C(p, \theta, T,R) \left[ \mathcal{M}_1^p + \mathcal{M}_2^p e^{apT} \left(2 +  \frac{c(\sigma)}{L(\sigma)}\right)^p\right].
		\]
	\end{proof}
	
	\begin{remark}
		The H\"older continuity in time proven above in Proposition \ref{P:Holder} recovers the result proven by Foondun and Nane in \cite[Theorem 1.9]{FN2016} if we restrict ourself to the case where $\beta \in (0,1)$ and $\gamma = 1-\beta$. 
	\end{remark}
	
	\begin{NickQuestion}
	\label{Q:HolderGeneralInitial}
		Suppose we consider bounded initial conditions $u_0$ and $v_0$. Then 
		\begin{align*}
		J_0(t,x) &=\begin{cases} 
		[Z(t,\cdot)*u_0](x) & \beta \in (0,1] \\
		[Z(t,\cdot)*v_0](x) + \frac{\partial}{\partial t}[Z(t,\cdot)*u_0](x) & \beta \in (1,2)
		\end{cases} \\
		& =: \begin{cases} 
		[Z(t,\cdot)*u_0](x) & \beta \in (0,1] \\
		[Z(t,\cdot)*v_0](x) + [Z^*(t,\cdot)*u_0](x) & \beta \in (1,2)
		\end{cases}.
		\end{align*} 
		We notice that 
		\begin{align*}
		J_0(t,x) - J_0(s,z)
		&= \begin{cases} 
		[Z(t,\cdot)*u_0](x) - 	[Z(s,\cdot)*u_0](z) & \beta \in (0,1] \\ \\
		[Z(t,\cdot)*v_0](x) - [Z(s,\cdot)*v_0](z) \\
		\qquad + [Z^*(s,\cdot)*u_0](z) - [Z^*(t,\cdot)*u_0](x) & \beta \in (1,2)
		\end{cases}.
		\end{align*}
		Lets just examine the $\beta \in (1,2)$ case since that case deals with both $Z$ and $Z^*$. We have that 
		\begin{align*}
		|J_0(t,x) - J_0(s,z)| & \le \int_{\R} |Z(t,x-y) - Z(s,z-y)| |v_0(y)| \ud y + \int_{\R} |Z^*(t,x-y) - Z^*(s,z-y)| |u_0(y)| \ud y.
		\end{align*}
		For simplicity lets look at the time and space increments separately. For time we have that 
		\begin{align*}
		|J_0(t,x) - J_0(s,x)| & \le \int_{\R} |Z(t,x-y) - Z(s,x-y)| |v_0(y)| \ud y + \int_{\R} |Z^*(t,x-y) - Z^*(s,x-y)| |u_0(y)| \ud y.
		\end{align*}
		Foondun and Nane \cite[Proof of Theorem 1.13]{FN2016} stop right and claim that the above integrals are smooth functions in $t$. Then it seems that they conclude that they are $\eta$-H\"older continuous in $t$ with $\eta=\frac{(1-\beta \gamma/\alpha)}{2}$ where $\gamma$ is the Reisz Kernel exponent. In other words they say that 
		\[
		|J_0(t,x) - J_0(s,x)| \le C |t-s|^\eta.
		\]
		
		\textcolor{blue}{Are smooth functions also $\eta$-H\"older continuous for any $\eta >0$?} If we can argue as in \cite{FN2016}, then we can relax our initial conditions to be bounded measurable functions $u_0$ and $v_0$. This is also an improvement from Sanz-Sole and Millet \cite{SoleMillet2021} since they assume the initial conditions are continuous. 
	\end{NickQuestion}

	\begin{NickQuestion}[Continue of Question \ref{Q:HolderGeneralInitial}]
Suppose that $\beta \in (0,1]$.

\noindent \textbf{Time Increment}: Suppose $0 \le s < t$. Then,
	\begin{align}
		|J(t,x) - J(s,x)| &\le \Norm{u_0}_{\infty} \int_{\R^d} | Z(t,x-y) - Z(s, x-y) | \ud y \nonumber \\
		&= \Norm{u_0}_{\infty} |t-s| \int_{\R^d} \left| \frac{\partial}{\partial t} Z(\xi(s,t),x-y) \right| \ud y \label{E:DZnorm}
	\end{align}
where $s \le \xi(s,t) \le t$. Now we will calculate $\frac{\partial}{\partial t} Z(\xi(s,t),x-y)$.
	\begin{align*}
		\frac{\partial}{\partial t} Z(\xi(s,t),x) &= 	\frac{\partial}{\partial t} \left( \pi^{-d/2} t^{\lceil \beta \rceil - 1} |x|^{-d} H_{2,3}^{2,1}\left( \frac{|x|^\alpha}{2^{\alpha -1}\nu t^\beta} \bigg\vert \begin{matrix} (1,1) & (\lceil \beta \rceil,\beta) &
		\\ (d/2,\alpha/2) & (1,1) & (1,\alpha/2) \end{matrix} \right)  \right) \\
		&= \frac{\partial}{\partial t} \left( \pi^{-d/2} |x|^{-d} H_{2,3}^{2,1}\left( \frac{|x|^\alpha}{2^{\alpha -1}\nu t^\beta} \bigg\vert \begin{matrix} (1,1) & (1,\beta) &
		\\ (d/2,\alpha/2) & (1,1) & (1,\alpha/2) \end{matrix} \right)  \right) \\
		&= \pi^{-d/2} |x|^{-d} \frac{\partial}{\partial t} \left( H_{2,3}^{2,1}\left( \frac{1}{(ct)^\beta} \bigg\vert \begin{matrix} (1,1) & (1,\beta) &
		\\ (d/2,\alpha/2) & (1,1) & (1,\alpha/2) \end{matrix} \right)  \right)
	\end{align*}
where $ c = (2^{\alpha-1}\nu / |x|^{-\alpha})^{1/\beta}$. We now apply \cite[(2.2.6)]{Kilbas} to see that,
	\begin{align*}
	\pi^{-d/2} |x|^{-d} \frac{\partial}{\partial t}  H_{2,3}^{2,1} & \left( \frac{1}{(ct)^\beta} \bigg\vert \begin{matrix} (1,1) & (1,\beta) &
		\\ (d/2,\alpha/2) & (1,1) & (1,\alpha/2) \end{matrix} \right) \\  
		&= \frac{\pi^{-d/2}|x|^{-d} }{t}  H_{3,4}^{3,1}\left( \frac{1}{(ct)^\beta} \bigg\vert \begin{matrix} (1,1) & (1,\beta) & (0, \beta) &
			\\(1,\beta) & (d/2,\alpha/2) & (1,1) & (1,\alpha/2) \end{matrix} \right) \\
		&= \frac{\pi^{-d/2}|x|^{-d} }{t}  H_{3,4}^{3,1}\left( \frac{1}{(ct)^\beta} \bigg\vert \begin{matrix} (1,1) & (0, \beta) & (1,\beta) &
			\\(1,\beta) & (d/2,\alpha/2) & (1,1) & (1,\alpha/2) \end{matrix} \right) \\
		& = \frac{\pi^{-d/2}|x|^{-d} }{t}  H_{2,3}^{2,1}\left( \frac{1}{(ct)^\beta} \bigg\vert \begin{matrix} (1,1) & (0, \beta) &
		\\ (d/2,\alpha/2) & (1,1) & (1,\alpha/2) \end{matrix} \right),
	\end{align*}
where in the last line we used \cite[(2.1.2)]{Kilbas}. Replacing $c$ above gives us that 
	\begin{align*}
		\frac{\partial}{\partial t} Z(\xi(s,t),x) &= \frac{\pi^{-d/2}|x|^{-d} }{t}  H_{2,3}^{2,1}\left( \frac{1}{(ct)^\beta} \bigg\vert \begin{matrix} (1,1) & (0, \beta) &
		\\ (d/2,\alpha/2) & (1,1) & (1,\alpha/2) \end{matrix} \right) \\ 
		& = \frac{\pi^{-d/2}|x|^{-d} }{t}  H_{2,3}^{2,1}\left( \frac{|x|^{\alpha}}{2^{\alpha -1} \nu t^\beta} \bigg\vert \begin{matrix} (1,1) & (0, \beta) &
		\\ (d/2,\alpha/2) & (1,1) & (1,\alpha/2) \end{matrix} \right).
	\end{align*}
We need to calculate \eqref{E:DZnorm}. In order to do this, we can calculate $\mathcal{F} \left( \frac{\partial}{\partial t} Z(\xi(s,t),\cdot) \right)(0)$. However if we examine the proof of \cite[Lemma 4.2]{CHHH16}, then is appears that 
	\[
		\mathcal{F} \left( \frac{\partial}{\partial t} Z(\xi(s,t),\cdot) \right)(\xi) = 
			\frac{1}{t} H_{1,2}^{1,1}\left( 2^{-1} \nu t^\beta |\xi|^{\alpha} \bigg\vert \begin{matrix} (0,1) & 
		\\ (0,1) & (1, \beta) \end{matrix} \right).
	\]
This can no longer be written as a Mittag-Leffler function since by \cite[(A.19)]{CHHH16} we would have 
	\[
		\frac{1}{t} H_{1,2}^{1,1}\left( 2^{-1} \nu t^\beta |\xi|^{\alpha} \bigg\vert \begin{matrix} (0,1) & 
		\\ (0,1) & (1, \beta) \end{matrix} \right) = \frac{1}{t} E_{\beta,0}\left( - 2^{-1} \nu t^\beta |\xi|^{\alpha} \right),
	\]
and $E_{\alpha, \beta}$ is not defined for $\beta = 0$. So I am unsure of how to calculate \eqref{E:DZnorm} in this case. I need to show that 
	\[
		\mathcal{F} \left( \frac{\partial}{\partial t} Z(\xi(s,t),\cdot) \right)(0) = 
		\frac{1}{t} H_{1,2}^{1,1}\left( 0 \bigg\vert \begin{matrix} (0,1) & 
		\\ (0,1) & (1, \beta) \end{matrix} \right) < \infty.
	\]
	\end{NickQuestion}
	
	\subsection{Existence of a Global Solution with Superlinear Diffusion Coefficient}
	
	We now impose the following condition on $\sigma$: As $|x|,|z| \to \infty$,
	\begin{equation}
	\label{C:SigSupLin}
	|\sigma(x)-\sigma(z)| \le \sigma_2|x-z|[\ln_+(|x-z|)]^{\delta},
	\end{equation}
	where $\sigma_2 , \delta >0$ and $\ln_+(z) = \ln(z \vee e)$ for $z >0$. Further, this implies that as $|x| \to \infty$,
	\begin{equation}
	|\sigma(x)| \le \sigma_1 + \sigma_2|x|[\ln(|x|)]^{\delta},
	\end{equation}
	where $\sigma_1 = c(\sigma) = \sigma(0)$.
	
	\begin{remark}
		Because of the absence of the drift coefficient in \eqref{E:SPDE}, we may drop \cite[Assumption C1]{SoleMillet2021}.
	\end{remark}
	
	\begin{theorem}
		Consider the initial conditions $u_0=1$ and $v_0=0$ and suppose that $\sigma$ satisfies \eqref{C:SigSupLin}. Then for any $M>0$, there exists a random field solution to \eqref{E:SPDE} in $[-M,M]$ denoted as $(u(t,x): (t,x) \in [0,T]\times[-M,M])$. This solution is unique and satisfies 
		\begin{equation}
		\sup_{t\in[0,T],\; x\in[-R,R]} |u(t,x)| < \infty, a.s.
		\end{equation}
	\end{theorem}
	
	\begin{proof}
		We first solve a {\it truncated} version of \eqref{E:SPDE}. By this, we mean $\eqref{E:SPDE}$ with $\sigma$ replaced by $\sigma_N$ where 
		\[
		\sigma_N(x) = \sigma(x)\textbf{1}_{\{|x| \le N\}} + \sigma(N)\textbf{1}_{\{x  >N\}} + \sigma(-N)\textbf{1}_{\{x \le -N\}}.
		\]
		Note that since $\sigma$ satisfies \eqref{C:SigSupLin}, then $\sigma_N$ is Lipschitz continuous with Lipschitz constant $L(\sigma_N) := \sigma_2 \ln(2N)^{\delta}$. In other words, 
		\[
		|\sigma_N(x)| \le \sigma_1 + \sigma_2 \ln(2N)^{\delta}|x|.
		\]
		Hence we can apply Theorem \ref{T:SolveabilityGenLip} and see that there exists an unique solution of the truncated version of \eqref{E:SPDE}, which we denote by $u_N := \left\{ u_N(t,x) : (t,x) \in [0,T] \times \R \right\}$. In addition, Proposition \ref{P:Holder} implies that the solution has a version, still denoted by $u_N$, which is $\eta_1$-H\"older continuous in time and $\eta_2$-H\"older continuous in space  with $0<\eta_1 < q/2$ and $0<\eta_2 < \theta/2$.
		
		We will now apply Proposition \ref{P:Holder} to find an upper bound on the $p$-norm of $u_N$. For this, consider 
		\[
		a > \left( 8 L^2(\sigma_N) K^2 \right)^{1/\rho} \quad \text{and} \quad p \in \left[ 2, \frac{a^{ \rho}}{4L^2(\sigma_N)K^2} \right].
		\] 
		Then Proposition \ref{P:Holder} implies that 
		\begin{equation}
		\label{E:u_NMomBdd}
		\mathbb{E}\left(\sup_{t\in[0,T],\; x\in[-R,R]} |u_N(t,x)|^p\right) \le 2^{p-1} + C(p, \theta, T,R)\left[\mathcal{M}_1^p + \mathcal{M}_2^p(N) \exp(apT)\left(2+ \frac{\sigma_1}{L(\sigma_N)}\right)\right]
		\end{equation}
		where we recall that 
		\begin{equation}
		\label{D:M1&2forSig_N}
		\mathcal{M}_1 = \sqrt{p}\sigma_1 \quad \text{and} \quad \mathcal{M}_2(N) = \sqrt{pL(\sigma_N)} = \sqrt{p\sigma_2\ln(2N)^\delta}.
		\end{equation}
		We now use the above to prove the existence and uniqueness of \eqref{E:SPDE} with superlinear $\sigma$ satisfying \eqref{C:SigSupLin}. For any $T>0$ and $N \ge 2$ with $N\in \mathbb{N}$, we define the following stopping time:
		\begin{equation}
		\tau_N := \inf \left\{ t>0 : \sup_{|x| \le R} |u_N(t,x)| \ge N \right\} \wedge T.
		\end{equation}
		The uniqueness of the solution and the \textcolor{blue}{local property of the stochastic integral} imply that on $\{t < \tau_N\}$, $u_N(t,x) = u_{N+k}(t,x)$ for any $k \in \mathbb{N}$. Hence, $(\tau_N)_{N \ge 2}$ is increasing and bounded above by $T$. 
		
		We now momentarily assume that $\sup_N \tau_N = T$ which would imply that $\{t < \tau_N \} \uparrow \Omega$. On each $\{ t < \tau_N \}$, define $(u(t,x) : (t,x) \in [0,T) \times \R)$ by $u(t,x) = u_N(t,x)$ and hence $u(t,x) = u_{N+k}(t,x)$ for any $k\in \mathbb{N}$. This implies that on $\{ t < \tau_N \}$ 
		\[
		u(t,x) = 1 + \int_0^t \ud s \int_{\R} \ud y \; Y(t-s,x-y) \sigma_N(u(s,y)) W(\ud s, \ud y).
		\]
		However, on $\{ t < \tau_N \}$, we have that $\sigma_N(u_N(t,x)) = \sigma(u(t,x))$ and so $u$ satisfies
		\[
		u(t,x) = 1 + \int_0^t \ud s \int_{\R} \ud y \; Y(t-s,x-y) \sigma(u(s,y)) W(\ud s, \ud y), \quad \{ t < \tau_N \}.
		\]
		Lastly, since $\{ t < \tau_N \} \uparrow \Omega$, we conclude that 
		\begin{equation}
		\label{E:SupLinMildSol}
		u(t,x) = 1 + \int_0^t \ud s \int_{\R} \ud y \; Y(t-s,x-y) \sigma(u(s,y)) W(\ud s, \ud y), \quad (t,x) \in [0,T) \times \R.
		\end{equation}
		\textcolor{blue}{It is mentioned that the stochastic integrals are no longer defined in $L^2(\Omega)$. The extension of the integral is found in Dalang and Sanz-Sole's textbook that is not yet published.}
		
		We now focus on proving that $\sup_N \tau_N = T$, or equivalently that $ \mathbb{P}\left(\tau_N <T\right) \to 0$ as $N \to \infty$. To do this, note that by Chebychev's inequality
		\begin{align*}
		\mathbb{P}\left( \tau_N < T \right) &\le \mathbb{P}\left(\sup_{t\in[0,T],\; x\in[-R,R]} |u_N(t,x)| \ge N \right) \\
		& \le N^{-p}\mathbb{E}\left( \sup_{t\in[0,T],\; x\in[-R,R]} |u_N(t,x)|^p \right).
		\end{align*}
		An application of \eqref{E:u_NMomBdd} gives 
		\begin{align*}
		N^{-p}\mathbb{E}\left( \sup_{t\in[0,T],\; x\in[-R,R]} |u_N(t,x)|^p \right) &\le N^{-p}\left(  2^{p-1} + C(p, \theta, T,R) \left[ \mathcal{M}_1^p + \mathcal{M}_2(N)^p \exp\{apT\} \left(2 +  \frac{\sigma_1}{L(\sigma_N)}\right)^p\right]\right) \\
		&= \frac{2^{p-1} + C(p, \theta, T,R)\mathcal{M}_1^p }{N^p} + \frac{C(p,T,R) \mathcal{M}_2(N)^p \exp\{apT\} \left(2 +  \frac{\sigma_1}{L(\sigma_N)}\right)^p }{ N^p }.
		\end{align*}
		The first term clearly converges to $0$ as $N\to \infty$. As for the second term, we expand using \eqref{D:M1&2forSig_N} to see that 
		\begin{align*}
		\frac{\mathcal{M}_2(N)^p \exp\{apT\} \left(2 +  \frac{\sigma_1}{L(\sigma_N)}\right)^p }{ N^p } &= \frac{\left(\sqrt{p\sigma_2\ln(2N)^\delta}\right)^p \exp\{apT\} \left(2 +  \frac{\sigma_1}{ \sqrt{p\sigma_2\ln(2N)^\delta} }\right)^p }{ N^p }  \\
		&\sim \frac{ \ln(2N)^{p\delta/2} }{ N^p } \to \infty,
		\end{align*}
		where above we use the symbol $\sim$ to mean asymptotically equivalent as $N \to \infty$. This completes the proof of the Theorem. 
	\end{proof}
	
	\begin{remark}
		A potential project would be to show that the solution to \eqref{E:SPDE} with superlinear coefficient is Holder continuous. This requires to one to extend the Kolmogorov theorem to the type of stochastic integral used to define the mild solution in \eqref{E:SupLinMildSol}. However these stochastic integrals are no longer defined in $L^2(\Omega)$. Maybe a completely new theory needs to be developed to show continuity. \cite[Theorem 5.5.5]{Kuo} may be helpful? 
	\end{remark}
	
	\section{Spatially Colored Noise Driven SPDE}
In this section we will study \eqref{E:SPDE} driven by a Gaussian noise that is white in time and colored in space. In other words, we start with a non-negative and nonnegative definite tempered distribution on $\mathscr{S}'(\R^d)$ which we denote as $\Lambda$. Then the Bochner-Schwartz theorem says that $\Lambda$ is the Fourier transform of a non-negative, tempered and symmetric measure which we denote as $\mu$. $\mu$ is referred to as the {\it spectral measure}. Because of this, the following functional is nonnegative definite and symmetric on $C_0(\R_+ \times \R^d)$:
	\begin{equation}
	\label{E:CovNoise}
		C(\phi, \psi) := \int_0^\infty \ud t \int_{\R^{d}}  \Lambda(\ud x) \; (\phi(t,\cdot)*\tilde{\psi}(t,\cdot))(x), \quad \tilde{\psi}(x) = \psi(-x).
	\end{equation}
Thus we may identify $C$ as the covariance functional of some Gaussian process. This leads to the following definition.
	\begin{definition}
A zero-mean Gaussian processes, $F := \{ F(\phi) : \phi \in C_0(\R_+ \times \R^d) \}$ is said to be white in time and homogeneously colored in space if $\mathbb{E}(F(\phi)F(\psi)) = C(\phi,\psi)$. 
	\end{definition}

	\begin{assumption}
	\label{A:NoiseAssum1}
For the remainder of this paper, we will only consider the Gaussian noises that satisfy the following three conditions:
		\begin{enumerate}
			\item the spectral measure $\mu$ is such that $\int_{\R^d} \frac{\mu(\ud \zeta)}{1 + |\zeta|^\alpha} < \infty$,
			\item $\Lambda$ is absolutely continuous with respect to the Lebesgue measure and therefore has a density, $\Lambda(\ud x) = f(x) \ud x$. Furthermore, when this is the case, we may often write $\mu(\ud \zeta) = \varphi(\zeta) \ud \zeta$ where we have the relation $\mathcal{F}\varphi = f$,
			\item the correlation functions satisfies for some $\lambda \in (0,d)$ and for all $c>0$, $f(cx) = c^{-\lambda}f(x)$, which in return implies that $\varphi(cx) = c^{-(d- \lambda)}\varphi(x)$.
		\end{enumerate}
	\end{assumption}

\noindent Under this assumption, we may rewrite \eqref{E:CovNoise} as 
	\[
		C(\phi,\psi) = \int_0^\infty \ud t \int_{\R^{2d}} \ud x \ud y \; \phi(t,x) f(x-y) \psi(t,y)
	\]
  and by applying Placherel's identity, this further implies that 
  	\[
  		\int_0^\infty \ud t \int_{\R^{2d}} \ud x \ud y \; \phi(x) f(x-y) \psi(y) = (2 \pi)^{-d} \int_0^\infty \ud t \int_{\R^{2d}} \mu(\ud \zeta) \mathcal{F}(\phi(t,\cdot))(\zeta) \overline{\mathcal{F}(\psi(t,\cdot))(\zeta)}.
  	\]
We now prove a lemma which will be used to prove the existence and uniqueness of the solution. It is proven by Foondun and Nane \cite{FN2016} in the case that 		
\begin{equation}
\label{A:Y=GAssumptions}
\beta \in (0,1) \quad \text{and} \quad \gamma = 1 - \beta.
\end{equation}		

  	\begin{lemma}
  	\label{L:YfYbdd}
  	Suppose Assumption \ref{A:NoiseAssum1} holds. Then there exists a constant $C(\beta,\gamma) =: C >0$ such that for any $x,y \in \R^d$ we have that 
  	\begin{equation}
  	\label{E:YYfBdd}
  	\int_{\R^{2d}} Y(t,x-z)Y(t,y-w)f(w-z) \; \ud z \ud w \le C \left( \frac{\nu}{2} \right)^{-\lambda / \alpha} t^{2(\beta + \gamma -1)-\beta \lambda / \alpha} \int_{\R^{d}}  \frac{ \mu(\ud \zeta) }{\left(1+ |\xi|^\alpha\right)^2 } < \infty
  	\end{equation}
  \end{lemma}
  \begin{proof}
 Through a change of variables, we may write the left side of \eqref{E:YYfBdd} as the following:
  	\[
  	\int_{\R^{2d}} Y(t,x-z)Y(t,y-w)f(w-z) \; \ud z \ud w = \int_{\R^{d}} Y\left(t,u-(x-y)\right)(Y(t,\cdot)*f)(u) \; \ud u,
  	\]
 where $*$ denotes the spatial convolution. Now an application of Pareseval's identity gives us that 
  	\[
  	\int_{\R^{d}} Y\left(t,u-(x-y)\right)(Y(t,\cdot)*f)(u) \; \ud u = \frac{1}{(2\pi)^d} 	\int_{\R^{d}} \mathcal{F}Y\left(t,\cdot-(x-y)\right)(\zeta) \mathcal{F}(Y(t,\cdot)*f)(\zeta) \ud \zeta.
  	\]
 Applying the translation property of the Fourier transform and \cite[Theorem 4.1]{CHN19} gives us that 
  	\begin{align*}
  	\frac{1}{(2\pi)^d} 	\int_{\R^{d}} \mathcal{F}Y\left(t,\cdot-(x-y)\right)(\zeta) & \mathcal{F}(Y(t,\cdot)*f)(\zeta) \ud \xi \\
  	& = \frac{ t^{2(\beta + \gamma -1)}}{ (2\pi)^d } \int_{\R^{d}} \exp(-i(x-y)\zeta) E_{\beta,\beta+\gamma}\left(-\frac{\nu t^\beta |\zeta|^\alpha}{2}\right)^2 \mu(\ud \zeta).
  	\end{align*}
 Hence by taking the absoulte value on both sides and recalling the positivity of $f$, $Y$ and the Mittag-Leffler function we get that 
  	\begin{align*}
  			\int_{\R^{2d}} Y(t,x-z)Y(t,y-w)f(w-z) \; \ud z \ud w \le \frac{ t^{2(\beta + \gamma -1)} }{ (2\pi)^d } \int_{\R^{d}}  E_{\beta,\beta+\gamma}\left(-\frac{\nu t^\beta |\zeta|^\alpha}{2}\right)^2 \mu(\ud \zeta).
  	\end{align*}
Now, by applying \cite[Theorem 1.6]{Podlubny99}, we see that there exists a $C>0$ depending only on $\beta$ and $\gamma$ such that 
	\[
		\frac{ t^{2(\beta + \gamma -1)}}{ (2\pi)^d } \int_{\R^{d}}  E_{\beta,\beta+\gamma}\left(-\frac{\nu t^\beta |\xi|^\alpha}{2}\right)^2 \mu(\ud \zeta).\le 
  		\frac{C  t^{2(\beta + \gamma -1)}}{ (2\pi)^d } \int_{\R^{d}}  \frac{ \mu(\ud \zeta) }{\left(1+ \frac{\nu}{2} t^\beta |\zeta|^\alpha \right)^2 } .
 	\]
 Another change of variables implies that 
  	\[
  	\frac{C  t^{2(\beta + \gamma -1)}}{ (2\pi)^d } \int_{\R^{d}}  \frac{ \mu(\ud \zeta) }{\left(1+ \frac{\nu}{2} t^\beta |\zeta|^\alpha \right)^2 }  \le 
  	\frac{C }{ (2\pi)^d } \left(\frac{\nu}{2}\right)^{-\lambda/\alpha}  t^{2(\beta + \gamma -1)-\beta \lambda/\alpha} \int_{\R^{d}}  \frac{ \mu(\ud \zeta) }{\left(1+ |\xi|^\alpha\right)^2 } .
  	\]
  \textcolor{blue}{? Recall that $\varphi(c \zeta) = c^{-d+\lambda}\varphi(\zeta)$ and so when doing the substitution $\eta = \left(\frac{\nu}{2}t^{\beta} \right)^{1/\alpha} \zeta$ then we should get 
  		\begin{align*}
  			\frac{C  t^{2(\beta + \gamma -1)}}{ (2\pi)^d } \int_{\R^{d}}  \frac{ \mu(\ud \zeta) }{\left(1+ \frac{\nu}{2} t^\beta |\zeta|^\alpha \right)^2 } &= \frac{C  t^{2(\beta + \gamma -1)}}{ (2\pi)^d } \int_{\R^{d}}  \frac{ \varphi\left( \left(\frac{\nu}{2}t^{\beta} \right)^{-1/\alpha}\eta\right)\ud \zeta }{\left(1+  \left|\eta \right|^\alpha \right)^2 \left(\frac{\nu}{2}t^{\beta} \right)^{d/\alpha} } \\
  			&= \frac{C  t^{2(\beta + \gamma -1)}}{ (2\pi)^d } \int_{\R^{d}}  \frac{ \left(\frac{\nu}{2}t^{\beta} \right)^{(d-\lambda)/\alpha} \varphi\left( \eta\right)\ud \zeta }{\left(1+  \left|\eta \right|^\alpha \right)^2 \left(\frac{\nu}{2}t^{\beta} \right)^{d/\alpha} } \\
  			&=\left(\frac{\nu}{2}\right)^{\frac{-\lambda }{\alpha}}\frac{C  t^{2(\beta + \gamma -1)-\frac{\lambda \beta}{\alpha}}}{ (2\pi)^d } \int_{\R^{d}}  \frac{ \varphi\left( \eta\right)\ud \zeta }{\left(1+  \left|\eta \right|^\alpha \right)^2 } 
  		\end{align*}
   }
The finiteness of the above integral follows from Assumption \ref{A:NoiseAssum1}.
  \end{proof}	
  	
  \begin{NickQuestion}
 Lemma \ref{L:YfYbdd} is a problem. Later on we will need to assume that 
 	\begin{equation}
 	\label{A:Problem}
 		2(\beta + \gamma -1)-\beta \lambda/\alpha > 0
 	\end{equation}
 which is stronger than the original assumption that 
 	\[
 		2(\beta + \gamma) -1 -\beta \lambda/\alpha.
 	\]
 \eqref{A:Problem} excludes the SHE and the case when $\gamma = \lceil \beta \rceil -1$. 
  \end{NickQuestion}	
  	
	\begin{example}
	\label{E:RieszKernel}
Consider the function $f(x)$, which is the well-known Riesz kernel with parameter $\lambda \in (0,d)$:
		\[
			f(x) := \frac{1}{|x|^\lambda} \quad \text{and} \quad \varphi(x):= \mathcal{F}f(\xi) = \frac{K_f}{|\xi|^{d-\lambda}},
		\]
	where
		\[
			K_f     = \pi^{-d/2}2^{-\lambda} \frac{\Gamma((d-\lambda)/2)}{\Gamma(\lambda/2)}.
		\]
Considering Assumption \ref{A:NoiseAssum1}, if we assume $\lambda < \min\{d,2\}$, then we may consider our Gaussian processes to have its correlation function given to be the Riesz kernel, $f(x)=|x|^{-\lambda}$. 
	\end{example} 

	\subsection{Lipschitz Diffusion Coefficient}
	We first examine \eqref{E:SPDE} with Lipschitz diffusion coefficient . As we will see below in Theorem \ref{T:ExisAndUniq}, to establish the existence and uniqueness of a random field solution to \eqref{E:SPDE}, we will need to assume Dalang's condition:
	\begin{equation}
	\label{C:DalangCl}
	\int_0^t \ud s \int_{\R^d} \ud y_1 \ud y_2 Y(s,y_1) f(y_1-y_2) Y(s,y_2 )< \infty, \quad \text{for all } t>0,
	\end{equation} 
	which is equivalent to
	\begin{equation}
	\label{E:DalangEquiv2Cl}
	\rho(\lambda) > 0 \quad \text{and} \quad \lambda < 2\alpha,
	\end{equation}
	where we recall from \eqref{D:theta} that
	\begin{equation}
	\rho(x) := 2\beta + 2\gamma -1 -\beta x / \alpha.
	\end{equation}
	We assume that there exists some $L_\sigma > 0$ such that 
	\[
	|\sigma(x) - \sigma(y)| \le L_{\sigma} |x-y|.
	\]


	We can now prove the existence and uniqueness of the solution to \eqref{E:SPDE} driven by the noise $\dot{F}$. 
	
	\begin{theorem}
		\label{T:ExisAndUniq}
		Suppose that $\lambda \beta / \alpha - 2(\beta + \gamma -1) < 1$ and in addition, suppose that $\lambda < 2 \alpha$. Then there exists a unique $L^2(\Omega)$-solution to \eqref{E:SPDE} driven by $\dot{F}$ such that 
		\[
		u(t,x) = J_0(t,x) + I(t,x), \quad a.s.
		\]
	\end{theorem} 
	\begin{proof}
		The proof follows the strategy laid out in Walsh \cite{Walsh89}. We start by defining 
		\[
		u_0(t,x) := J_0(t,x) \quad \text{and} \quad u_n(t,x) := J_0(t,x) + \int_0^t \int_{\R^d} Y(t-s,x-y) \sigma(u_n(s,y)) F(\ud s, \ud y). 
		\]
		In addition, we define 
		\[
		D_n(t,x) := \mathbb{E}\left| u_{n+1}(t,x) - u_n(t,x) \right|^2 \quad \text{and} \quad H_n(t) := \sup_{x \in \R^d} D_n(t,x).
		\]
		Lastly, we define 
		\[
		\mathbf{\Sigma} (t,x,n) := | \sigma( u_n(t,x) ) - \sigma(u_{n-1} (t,x) ) |.
		\]
		With this notation we see that 
		\[
		D_n(t,x) = \int_0^t \ud s \int_{\R^d} \ud z \; \ud y \; Y(t-s,x-y) Y(t-s,x-z) f(y-z) \mathbb{E}\left(\mathbf{\Sigma}(s,y,n) \mathbf{\Sigma}(s,z,n)\right).
		\]
		Now note that by the Lipschitz property of $\sigma$, we have that  
		\begin{align*}
		\mathbb{E}\left(\mathbf{\Sigma}(s,y,n) \mathbf{\Sigma}(s,z,n)\right) &\le L^2_{\sigma} \mathbb{E} \left(|u_n(s,y)-u_{n-1}(s,y)||u_n(s,z)-u_{n-1}(s,z)|\right) \\
		& \le L^2_{\sigma} \Norm{ u_n(s,y)-u_{n-1}(s,y) } \Norm{u_n(s,z)-u_{n-1}(s,z)} \\ 
		& \le L^2_{\sigma} \sqrt{D_{n-1}(s,y) D_{n-1}(s,z)} \\
		& \le L^2_{\sigma} H_{n-1}(s).
		\end{align*}
		Thus, 
		\begin{align*}
		D_n(t,x) &\le L^2_{\sigma} \int_0^t H_{n-1}(s) \int_{\R^d} \ud z \; \ud y \; Y(t-s,x-y) Y(t-s,x-z) f(y-z) \\
		&\le C L^2_{\sigma} \left(\frac{\nu}{2}\right)^{-\lambda / \alpha} \int_0^t \ud s \; \frac{H_{n-1}(s)}{(t-s)^{\lambda \beta / \alpha - 2(\beta + \gamma -1)}},
		\end{align*}
		where the last inequality follows from Lemma \ref{L:YfYbdd}. In return, this implies that 
		\[
		H_n(t) \le C L^2_{\sigma} \left(\frac{\nu}{2}\right)^{-\lambda / \alpha} \int_0^t \ud s \; \frac{H_{n-1}(s)}{(t-s)^{\lambda \beta / \alpha - 2(\beta + \gamma -1)}}.
		\]
		Since we assume that $\lambda \beta / \alpha - 2(\beta + \gamma -1) < 1$, then the result follows from the discussion before the proof of \cite[Lemma 3.3]{Walsh89}. 
	\end{proof}
	
	\begin{remark}
		One should note that when considering the SHE ($\beta =1$, $\alpha=2$ and $\gamma=0$), the assumptions $\lambda \beta / \alpha - 2(\beta + \gamma -1) < 1$ and $\lambda < 2 \alpha$ reduce down to $\lambda<2$ and $\lambda <3$ which recovers the well-known condition that one must have $\lambda < \min\{2,d\}$.
	\end{remark}
	
	\begin{NickQuestion}
		In the existence and uniqueness proof in \cite[Theorem 1.10]{FN2016}, the following is shown 
		\[
		H_n(t) \le \text{const.} \int_0^t \frac{H_{n-1}(s)}{(t-s)^{\gamma \beta / \alpha}}.
		\]
		Note that in this formula, $\gamma$ is denoting the Riesz kernel parameter. I believe that we need to assume that 
		\[
		\gamma \beta / \alpha < 1
		\]
		however in this proof they say we must assume that $d\beta / \alpha < 1$. I think this is a typo. However, this assumption still is not stated in the statement of the theorem. Another issue is that they need to assume that 
		\[
		\gamma < \alpha
		\]
		in order to apply \cite[Lemma 2.7]{FN2016} and this assumption is not stated in the statement of the theorem. This assumption however would imply that $\gamma \beta / \alpha <1$ since in \cite{FN2016} $\beta \in (0,1)$.
		
		In addition, I believe that there is a typo in the statement of \cite[Lemma 3.3]{Walsh89}. I believe that in Lemma 3.3 {\it ibid.} we should assume that $a>-1$ and not $a >1$. This would be consistent with \cite[(3.14)]{Walsh89}.
	\end{NickQuestion}
	
	
	\textcolor{blue}{
		The following $L^2(\R^d)$ norms may be helpful. I may remove this. Not sure if I actually need it.
		\begin{lemma}
			\label{L:2NormZandZ*}
			Suppose that $d < 2 \alpha$ and that $\beta \in (0,2)$ and that $\gamma \le 0$. Then,
			\begin{equation}
			\label{E:2NormZ}
			\int_{\R^d} Z^2(t,x) \ud x = C_{Z} t^{\lceil \beta \rceil -1 - d\beta / \alpha}
			\end{equation}		
			where $C_{Z} = (2\pi)^{-d}(\nu/2)^{-d/\alpha}\int_{\R^d} E_{\beta,\lceil \beta \rceil}\left( -2^{-1}\nu t^\beta |\xi|^\alpha \right)^2 \ud \xi$. In addition, we have that
			\begin{equation}
			\label{E:2NormZ*}
			\int_{\R^d} {Z^*}^2(t,x) \ud x = C_{Z*} t^{-d\beta / \alpha}
			\end{equation}		
			where $C_{Z^*} = (2\pi)^{-d} (\nu/2)^{-d/\alpha} \int_{\R^d} E_{\beta,1}\left( -2^{-1}\nu t^\beta |\xi|^\alpha \right)^2 \ud \xi$.
		\end{lemma}
		\begin{proof}
			We apply Parseval's identity and \cite[Theorem 4.1]{CHN19} to see that 
			\begin{align*}
			\int_{\R^d} Z^2(t,x) \ud x &= \frac{t^{\lceil \beta \rceil -1}}{(2\pi)^d} \int_{\R^d} E_{\beta,\lceil \beta \rceil}\left( -2^{-1}\nu t^\beta |\xi|^\alpha \right)^2 \ud \xi \\
			&= \frac{t^{\lceil \beta \rceil -1 - d\beta / \alpha}}{(2\pi)^d} \left( \frac{\nu}{2} \right)^{-d/\alpha} \int_{\R^d} E_{\beta,\lceil \beta \rceil}\left( -|\xi|^\alpha \right)^2 \ud \xi,
			\end{align*}
			where in the last line we applied a change of variables. Note that by \cite[Theorem 1.6]{Podlubny99}, the integral above on the right hand side is finite when $2\alpha >d$. This proves \eqref{E:2NormZ}. Similarly we have 
			\begin{align*}
			\int_{\R^d} {Z^*}^2(t,x) \ud x &= \frac{1}{(2\pi)^d} \int_{\R^d} E_{\beta,1}\left( -2^{-1}\nu t^\beta |\xi|^\alpha \right)^2 \ud \xi \\
			&= \frac{t^{-d\beta /\alpha}}{(2\pi)^d} \left(\frac{\nu}{2}\right)^{-d/\alpha} \int_{\R^d} E_{\beta,1}\left( -2^{-1}\nu t^\beta |\xi|^\alpha \right)^2 \ud \xi,
			\end{align*}
			where again the integral above on the right hand side is finite when $d < 2\alpha$.
		\end{proof}
	}
	
	\subsection{Section 4.2 Millet and Sanz-Sol\'e \cite{SoleMillet2021}}
	
	\begin{proposition}{\normalfont\cite[Proposition 4.4]{SoleMillet2021}}
		Let $u_0$ and $v_0$ be Borel functions satisfying $\Norm{u_0}_\infty + \Norm{v_0}_\infty < \infty$. Suppose that $d < 2 \alpha$ and for this proof consider $\rho = \rho(\lambda)$ which is defined in \eqref{D:theta}. Then there exists a universal constant $K:=K(\alpha, \beta, \gamma,d, \lambda ) >0$ such that for any 
		\[
		a^\rho > \left( 8 L^2(\sigma)K^2\right).
		\]
		and 
		\[
		p \in \left[ 2, \frac{a^{\rho}}{4L^2(\sigma)K^2} \right].
		\]
		we have that 
		\begin{equation}
		\label{E:ColorNbdd}
		N_{a,p}(u)  \le 2 \mathcal{T}_0 +  \frac{c(\sigma)}{L(\sigma)}
		\end{equation}	
		where
		\begin{align*}
		\mathcal{T}_0 = \begin{cases} 
		(\lceil \beta \rceil - 1)^{\lceil \beta \rceil -1}(ea)^{1-\lceil \beta \rceil} \Norm{u_0}_\infty  & \beta \in (0,1]
		\\ (\lceil \beta \rceil - 1)^{\lceil \beta \rceil -1}(ea)^{1-\lceil \beta \rceil} \Norm{v_0}_\infty + \Norm{u_0}_\infty & \beta \in (1,2).
		\end{cases} 
		\end{align*}
		Moreover for any $T>0$,
		\begin{equation}
		\label{E:ColorPMomBdd}
		\sup_{x\in\R^d} \mathbb{E}(|u(t,x)|^p) \le \exp\left( a p t  \right) \left[ 2 \mathcal{T}_0 + \frac{ c(\sigma) }{L(\sigma)} \right]^p, \; t\in[0,T].
		\end{equation}
	\end{proposition}
	
	\begin{proof}
		We note that $\mathcal{N}_{a,p}(J)$ has the same upper bound as in Proposition \ref{P:N(u)Bdd} above. In other words,
		\[
		\mathcal{N}_{a,p}(J) \le \begin{cases}
		\Norm{u_0}_{\infty} & \beta \in (0,1] \\
		(ea)^{-1}\Norm{v_0}_{\infty} + \Norm{u_0}_{\infty} & \beta \in (1,2)
		\end{cases}.
		\]
		
		We know work to find an upper bound for $\mathcal{N}_{a,p}(I)$. The Burkholder-Davies-Gundy inequality and H\"older's inequality imply that 
		\begin{align*}
		\Norm{I(t,x)}_p^2 & \le 4p \int_0^t \ud s \iint_{\R^{2d}} \ud y_1 \ud y_2 Y(t-s,x-y_1)f(x-y)Y(t-s,x-y_2) \Norm{\sigma(u(s,y_1))}_p\Norm{\sigma(u(s,y_2))}_p.
		\end{align*} 
		We now apply the Lipschitz property to see that for $i\in\{1,2\}$,
		\begin{align*}
		\Norm{\sigma(u(s,y_i))}_p \le c(\sigma) + L(\sigma)\Norm{u(s,y_i)}_p.
		\end{align*}
		Therefore, we may say that 
		\begin{align*}
		\Norm{I(t,x)}_p^2 & \le 4p \int_0^t \ud s \iint_{\R^{2d}} \ud y_1 \ud y_2 Y(t-s,x-y_1)f(x-y)Y(t-s,x-y_2) \\
		& \hspace{10em} \times \prod_{i=1}^2 c(\sigma) + L(\sigma)\Norm{u(s,y_i)}_p \\
		& \le 4p \int_0^t \ud s \iint_{\R^{2d}} \ud y_1 \ud y_2 Y(t-s,x-y_1)f(x-y)Y(t-s,x-y_2) \left[c(\sigma) + L(\sigma)\exp(as)N_{a,p}(u)\right]^2 \\
		& \le 8p c(\sigma)^2 \int_0^t \ud s \iint_{\R^{2d}} \ud y_1 \ud y_2 Y(t-s,x-y_1)f(x-y)Y(t-s,x-y_2) \\
		& \qquad + 8p N_{a,p}(u)^2 L(\sigma)^2\int_0^t \ud s \exp(2as) \iint_{\R^{2d}} \ud y_1 \ud y_2 Y(t-s,x-y_1)f(x-y)Y(t-s,x-y_2).
		\end{align*}
		Now by applying Lemma \ref{L:intYfYbdd}, we may bound the $\iint_{\R^{2d}}\ud y_1 \ud y_2$ integrals by the quantity 
		\[
		K_1 (t-s)^{2(\beta + \gamma -1)-\beta \lambda /\alpha},
		\]
		where 
		\begin{equation*}
		K_1 : = C K_f \left(\frac{\nu}{2}\right)^{-\lambda/\alpha},
		\end{equation*}
		and $C$ is as in Lemma \ref{L:intYfYbdd}. In return, this gives us that 
		\begin{align*}
		\Norm{I(t,x)}_p^2 & \le  8p K_1 c(\sigma)^2 \int_0^t (t-s)^{2(\beta + \gamma -1)-\beta \lambda /\alpha} \ud s \\
		&\qquad + 8p K_1 N_{a,p}(u)^2 L(\sigma)^2\int_0^t \exp(2as) (t-s)^{2(\beta + \gamma -1)-\beta \lambda /\alpha} \ud s \\
		&= 8p K_1 c(\sigma)^2\frac{ t ^ \rho}{\rho} + 8p K_1 N_{a,p}(u)^2 L(\sigma)^2 \frac{\exp(2at) \left[ \Gamma(\rho) - \Gamma(\rho, 2at)\right]}{(2a)^{\rho} }  \\
		& =: \triangle_1(t) + \triangle_2(t),
		\end{align*}
		where $\Gamma(\cdot)$ denotes the Gamma function and $\Gamma(\cdot ,\circ)$ is the incomplete-upper Gamma function. Note that 
		\[
		\mathcal{N}_{a,p}(I)^2 \le \sup_{t \ge 0} \exp(-2at)\left( \triangle_1(t) + \triangle_2(t) \right).
		\]
		We start with bounding $\sup_{t \ge 0} \exp(-2at)\triangle_1(t)$. Recall the identities in \eqref{E:SUPidentities}. With this we see that 
		\[
		\sup_{t \ge 0} \exp(-2at) t^{\rho} = (\rho)^{\rho} (2ea)^{-\rho}.
		\]
		Hence, 
		\begin{align*}
		\sup_{t \ge 0} \exp(-2at)\triangle_1(t) & \le \frac{8p K_1 c(\sigma)^2}{\rho} \left(\frac{\rho}{2ea}\right)^{\rho}.
		\end{align*}
		Next we bound $\sup_{t \ge 0} \exp(-2at)\triangle_2(t)$. Note that for any $t\ge 0$,
		\begin{align*}
		\exp(-2at)\triangle_2(t) &= 8p K_1 N_{a,p}(u)^2 L(\sigma)^2 \frac{\left[ \Gamma(\rho) - \Gamma(\rho, 2at) \right]}{(2a)^\rho} \\
		& \le  8p K_1 N_{a,p}(u)^2 L(\sigma)^2 \frac{\Gamma(\rho)}{(2a)^{\rho}}.
		\end{align*}
		Putting this all together gives us that 
		\begin{align*}
		\mathcal{N}_{a,p}(I)^2 &\le \frac{8p K_1 c(\sigma)^2}{\rho} \left(\frac{\rho}{2ea}\right)^{\rho}  + \frac{8p K_1 \Gamma(\rho) L(\sigma)^2}{ (2a)^{\rho}}N_{a,p}(u)^2,
		\end{align*}
		and thus,
		\begin{equation*}
		\mathcal{N}_{a,p}(I)\le	2c(\sigma)\sqrt{\frac{2p K_1}{\rho}} \left(\frac{\rho}{2ea}\right)^{\rho/2}  + 2L(\sigma)\frac{\sqrt{2p K_1 \Gamma(\rho) }}{ (2a)^{\rho/2}}N_{a,p}(u).
		\end{equation*}
		Similar to the proof of Proposition \ref{P:N(u)Bdd} above, we define 
		\[
		K = K(\alpha,\beta,\gamma,d, \lambda):= \max \left\{2\sqrt{\frac{2 K_1}{\rho}} \left(\frac{\rho}{2e}\right)^{\rho/2},   2\frac{\sqrt{2 K_1 \Gamma(\rho) }}{ (2)^{\rho/2}} \right\}.
		\] 
		We now see that 
		\begin{equation}
		\mathcal{N}_{a,p}(I)\le K c(\sigma)\frac{\sqrt{p}}{ a^{\rho /2}} + K L(\sigma) \frac{\sqrt{p}}{ a^{\rho/2}} \mathcal{N}_{a,p}(u)
		\end{equation}
		This and the bound for $\mathcal{N}_{a,p}(J)$ found above implies that 
		\begin{align}
		\mathcal{N}_{a,p}(u) &\le \mathcal{N}_{a,p}(J) + \mathcal{N}_{a,p}(I) \nonumber \\
		& \le \begin{cases}
		\Norm{u_0}_{\infty} & \beta \in (0,1] \nonumber \\
		(ea)^{-1}\Norm{v_0}_{\infty} + \Norm{u_0}_{\infty} & \beta \in (1,2)
		\end{cases} \\
		& \qquad + K c(\sigma)\frac{\sqrt{p}}{ a^{\rho /2}} + K L(\sigma) \frac{\sqrt{p}}{ a^{\rho/2}} \mathcal{N}_{a,p}(u) \label{E:N(u)bddColored}.
		\end{align}
		Now consider 
		\[
		a^\rho > 8K^2L(\sigma)^2.
		\]
		This implies that the interval 
		\[
		\left[2, \frac{a^\rho}{4K^2L(\sigma)^2}  \right]
		\]
		is nonempty. Hence, for any $p$ in this interval we have that 
		\[
		\frac{KL(\sigma)\sqrt{p}}{a^{\rho/2}} \le \frac{1}{2}.
		\]
		This along with \eqref{E:N(u)bddColored} implies that 
		\begin{align*}
		\mathcal{N}_{a,p}(u) &\le 2\begin{cases}
		\Norm{u_0}_{\infty} & \beta \in (0,1] \nonumber \\
		(ea)^{-1}\Norm{v_0}_{\infty} + \Norm{u_0}_{\infty} & \beta \in (1,2)
		\end{cases} \\
		& \qquad + \frac{c(\sigma)}{L(\sigma)}.
		\end{align*}
		Its easy to see that \eqref{E:ColorPMomBdd} follows directly from \eqref{E:ColorNbdd}. 
	\end{proof}

We now state a lemma which will be used later on. 

	\begin{lemma}
	\label{L:Y_phiParseval}
Let $\phi$ and $\psi$ be bounded Borel functions. Then for $s,t > 0$ and $z \in \R^d$ we have,
	\begin{align}
		\iint_{\R^{2d}} \ud x \ud y \; \phi(x) Y(t,x) \psi(y) Y(s,y) f(x-y +z) = (2\pi)^{-d} \int_{\R^d} \mathcal{F}(\phi Y(t,\cdot))(\xi) \overline{\mathcal{F}(\psi Y(t,\cdot))(\xi)} e^{-iz \cdot \xi} \mu(\ud \xi)
	\end{align}
	\end{lemma}
	\begin{proof}
For notation, define $Y_\phi(t,x) = \phi(x)Y(t,x)$ and $Y_\psi(t,x) = \phi(x)Y(t,x)$. The boundedness and measureability of $\phi$ and $\psi$ guarantees the existence of $\mathcal{F}Y_\phi$ and $\mathcal{F}Y_\psi$. Thus we have,
	\begin{align*}
		\iint_{\R^{2d}} \ud x \ud y \; Y_{\phi}(t,x) Y_{\psi}(s,y) f(x-y +z) &= 	\int_{\R^{d}} \ud x \; Y_{\phi}(t,x) (Y_{\psi}(s,\cdot)*f(\cdot + z)) (x).
	\end{align*}		
Parseval's identity implies that 
	\begin{align*}
		\int_{\R^{d}} \ud x \; Y_{\phi}(t,x) (Y_{\psi}(s,\cdot)*f(\cdot + z)) (x) & = (2 \pi)^{-d}\int_{\R^{d}} \ud \zeta \; \mathcal{F}(Y_{\phi}(t,\cdot))(\zeta) \overline{\mathcal{F}((Y_{\psi}(s,\cdot))(\zeta) \mathcal{F}(f(\cdot + x))} \\
		&= (2 \pi)^{-d}\int_{\R^{d}} \mu(\ud \zeta) \; \mathcal{F}(Y_{\phi}(t,\cdot))(\zeta) \overline{\mathcal{F}((Y_{\psi}(s,\cdot))(\zeta)} e^{-ix\cdot\zeta},
	\end{align*}
where $\mu(\ud \xi) = \varphi(\xi) \ud \xi$. This completes the proof. 
	\end{proof}

	\begin{proposition}[\normalfont { \cite[Proposition 4.8]{SoleMillet2021} }]
	Statement...
	\end{proposition}

	\begin{proof}
The following proof follows closely the techniques presented in cite... . By the Burkholder Davies Gundy inequality, we have that 
	\begin{align*}
		\Norm{ I(t,x)-I(t,z) }_p^2 &\le 4p \Norm{ \int_0^t \ud s \int_{\R^{2d}} \ud y_1 \ud y_2 f(y_1-y_2) \prod_{i=1}^2 (Y(t-s,x-y_i) - Y(t-s, z-y_i)\sigma(u(s,y_i)) }_{p/2} \\
		& =: 4p \Norm{ Q }_{p/2} 
	\end{align*}
We now reorganize the integrand of $Q$ above by following \cite[Theorem 3.1]{HHN13}. By expanding the integrand of $Q$, we get that 
	\begin{align*}
		Q & =  \int_0^t \ud s \int_{\R^{2d}} \ud y_1 \ud y_2 \; Y(t-s,x-y_1) Y(t-s, x-y_2)f(y_1-y_2)\sigma(u(s,y_1))\sigma(u(s,y_2)) \\
		& \quad - \int_0^t \ud s \int_{\R^{2d}} \ud y_1 \ud y_2 \;  Y(t-s,x-y_1) Y(t-s, z-y_2)f(y_1-y_2)\sigma(u(s,y_1))\sigma(u(s,y_2)) \\
		& \quad - \int_0^t \ud s \int_{\R^{2d}} \ud y_1 \ud y_2 \;  Y(t-s,z-y_1) Y(t-s, x-y_2)f(y_1-y_2)\sigma(u(s,y_1))\sigma(u(s,y_2))  \\
		& \quad + \int_0^t \ud s \int_{\R^{2d}} \ud y_1 \ud y_2 \;  Y(t-s,z-y_1) Y(t-s, z-y_2)f(y_1-y_2)\sigma(u(s,y_1))\sigma(u(s,y_2)) \\
		& =: I_1 + I_2 + I_3 + I_4.
	\end{align*}
Now let $\xi = x-z$ and consider the following notation:
	\begin{align*}
	\Sigma_x(s,y) &= \sigma(u(s,x-y)), \quad \Sigma_{x,z}(s,y) = \sigma(u(s,x-y)) - \sigma(u(s,z-y)), \\
	h_1(s,y_1,y_2) &= f(y_1 - y_2) \Sigma_{x,z}(s,y_1)  \Sigma_{x,z}(s,y_2) \\
	h_2(s, y_1, y_2) &= [f(y_1 - y_2 + \xi) - f(y_1 - y_2)] \Sigma_x(s,y_2) \Sigma_{x,z}(s,y_1) \\ 
	h_3(s, y_1, y_2) &= h_2(s, y_2, y_1) \\
	h_4(s, y_1, y_2) &= [2f(y_1- y_2) - f(y_1 - y_2 + \xi) - f(y_1 - y_2 - \xi)]\Sigma_{x}(s,y) \Sigma_{x}(s,z).
	\end{align*}
With this notation along with a change of variables, we rewrite $I_1$ through $I_4$ as the following:
	\begin{align*}
		I_1 &= \int_0^t \ud s \int_{\R^{2d}} \ud y_1 \ud y_2 \; Y(t-s,y_1) Y(t-s, y_2)f(y_1-y_2)\Sigma_x(s,y_1)\Sigma_x(s,y_2) \\
		I_2 &= - \int_0^t \ud s \int_{\R^{2d}} \ud y_1 \ud y_2 \; Y(t-s,y_1) Y(t-s, y_2)f(y_1-y_2 - \xi)\Sigma_x(s,y_1)\Sigma_z(s,y_2) \\
		&= \int_0^t \ud s \int_{\R^{2d}} \ud y_1 \ud y_2 \; Y(t-s,y_1) Y(t-s, y_2)f(y_1-y_2 - \xi)\Sigma_x(s,y_1)\Sigma_{x,z}(s,y_2) \\
		& \quad - \int_0^t \ud s \int_{\R^{2d}} \ud y_1 \ud y_2 \; Y(t-s,y_1) Y(t-s, y_2)f(y_1-y_2 - \xi)\Sigma_x(s,y_1)\Sigma_{x}(s,y_2) \\
		I_3 &= \int_0^t \ud s \int_{\R^{2d}} \ud y_1 \ud y_2 \; Y(t-s,y_1) Y(t-s, y_2)f(y_1-y_2 + \xi)\Sigma_x(s,y_2)\Sigma_{x,z}(s,y_1) \\
		& \quad - \int_0^t \ud s \int_{\R^{2d}} \ud y_1 \ud y_2 \; Y(t-s,y_1) Y(t-s, y_2)f(y_1-y_2 + \xi)\Sigma_x(s,y_1)\Sigma_{x}(s,y_2) \\
		I_4 &= \int_0^t \ud s \int_{\R^{2d}} \ud y_1 \ud y_2 \; Y(t-s,y_1) Y(t-s, y_2)f(y_1-y_2)\Sigma_z(s,y_1)\Sigma_z(s,y_2) \\
		& = \int_0^t \ud s \int_{\R^{2d}} \ud y_1 \ud y_2 \; Y(t-s,y_1) Y(t-s, y_2)f(y_1-y_2)(\Sigma_{x,z}(s,y_1)-\Sigma_{x}(s,y_1)(\Sigma_{x,z}(s,y_2) - \Sigma_{x}(s,y_2)),
	\end{align*}
where above $I_3$ is derived in the same way $I_2$ was derived. Now if we expand $I_4$ and regroup all the terms in $I_1 + \cdots I_4$, we get that 
	\begin{equation}
		Q = \sum_{i=1}^4 Q_i
	\end{equation}
where
	\begin{equation}
	Q_i = \int_0^t \iint_{\R^{2d}} \ud y_1 \ud y_2 Y(t-s,y_1)Y(t-s,y_2)  h_i(s, y_1, y_2).
	\end{equation}
Thus 
	\begin{equation} \label{E:Ispaceinc}
		\Norm{ I(t,x)-I(t,z) }_p^2 \le 4p \sum_{i=1}^4 \Norm{Q_i}_{p/2}.
	\end{equation}
We now estimate each $\Norm{Q_i}_{p/2}$ separately.  

\textit{Upper bound for $\Norm{Q_1}_{p/2}$.} Using the Lipschitz property of $\sigma$, Minkowski's inequality and the triangle inequality gives is 
	\begin{align*}
		\Norm{Q_1}_{p/2} &\le L(\sigma)^2 \int_0^t \ud s \iint_{\R^{2d}} \ud y_1 \ud y_2 \; Y(t-s,y_1) Y(t-s, y_2) f(y_1 - y_2) \left[\sup_{|z_1-z_2|=\xi} \Norm{u(s,z_1)-u(s,z_2)}_p^2\right] \\
		& \le L(\sigma)^2C K_f \left(\frac{\nu}{2}\right)^{-\lambda/\alpha} \int_0^t \ud s \; (t-s)^{2(\beta + \gamma -1) - \beta \lambda / \alpha} \left[\sup_{|z_1-z_2|=\xi} \Norm{u(s,z_1)-u(s,z_2)}_p^2\right] \\
		& \le L(\sigma)^2C K_f \left(\frac{\nu}{2}\right)^{-\lambda/\alpha} (t)^{2(\beta + \gamma -1) - \beta \lambda / \alpha}\int_0^t \ud s \left[\sup_{|z_1-z_2|=\xi} \Norm{u(s,z_1)-u(s,z_2)}_p^2\right],
	\end{align*} 
where the second inequality above follows from Lemma \ref{L:intYfYbdd}. \textcolor{blue}{last line requires $2(\beta + \gamma -1) - \beta \lambda / \alpha > 0$. If we leave the $(t-s)^{2(\beta + \gamma -1) - \beta \lambda / \alpha}$ term in the integral then we may be able to assume the weaker condition that $2(\beta + \gamma -1) - \beta \lambda / \alpha > -1$, thus not excluding the SHE.}

To bound the remaining $Q_i's$, we will need to consider the following truncations which will then allow us to apply Lemma \ref{L:Y_phiParseval}. For $k \ge 1$ we define $\Sigma_{x}^k(s,y) = \Sigma_{x}(s,y) 1_{\{\Sigma_{x}(s,y) \le k \}}$ and $\Sigma_{x,z}^k(s,y) = \Sigma_{x,z}(s,y) 1_{\{\Sigma_{x,z}(s,y) \le k \}}$ and then we define $Q_i^k$ the same as $Q_i$ but with $h_i$ replaced with $h_i^k$ where $h_i^k$ is defined by replacing $\Sigma_{x}(s,y) $ and $\Sigma_{x,z}(s,y) $ with $\Sigma_{x}^k(s,y)$ and $\Sigma_{x,z}^k(s,y)$ respectively in $h_i$.

\textit{Upper Bound for $\Norm{Q_2^k}_{p/2}$.} We apply Lemma \ref{L:Y_phiParseval} with $\phi(y_2) = \Sigma_{x}^k(s,y_2)$ and $\psi(y_1) = \Sigma_{x,z}^k(s,y_1)$. Thus we have that up to a constant of $(2\pi)^{-d}$,
	\[
		Q_2^k = \int_0^t \ud s \int_{\R^d} \mu(\ud \zeta) \overline{ \mathcal{F}\left(\Sigma_{x}^k(s,\cdot)Y(t-s,\cdot)\right)(\zeta) }\mathcal{F}\left(\Sigma^k_{x,z} Y(t-s,\cdot)\right)(\zeta)[e^{-i\xi \cdot \zeta}-1],
	\]
where we recall that $\xi = x - z$. We now apply the the following to bounds: $|e^{-i\xi \cdot \zeta}-1| \le C |\xi|^{r} |\zeta|^{r}$ for any $r \in [0,1]$ and some $C>0$ and $\sqrt{ab} \le \frac{1}{2}(a+b)$ for and $a,b > 0$. Indeed, for some $r \in [0,1]$,
	\begin{align*}
		|Q_2^k| & \le C \int_0^t \ud s \int_{\R^d} \mu(\ud \zeta) \left| \mathcal{F}\left(\Sigma_{x}^k(s,\cdot)Y(t-s,\cdot)\right)(\zeta)\right| \left|\mathcal{F}\left(\Sigma^k_{x,z} Y(t-s,\cdot)\right)(\zeta)\right| |\xi|^{r} |\zeta|^{r} \\
		&\le \int_0^t \ud s |\xi|^{r} \left(\int_{\R^d} \mu(\ud \zeta) \left| \mathcal{F}\left(\Sigma_{x}^k(s,\cdot)Y(t-s,\cdot)\right)(\zeta)\right|^2 |\zeta|^{2r} \right)^{1/2} \\
		& \qquad \qquad \times \left( \int_{\R^d} \mu(\ud \zeta)\left|\mathcal{F}\left(\Sigma^k_{x,z} Y(t-s,\cdot)\right)(\zeta)\right|^2 \right)^{1/2} \\
		& \le \frac{1}{2} \int_0^t \ud s |\xi|^{2r} \int_{\R^d} \mu(\ud \zeta) \: |\zeta|^{2r} \left| \mathcal{F}\left(\Sigma_{x}^k(s,\cdot)Y(t-s,\cdot)\right)(\zeta)\right|^2  \\
		& \qquad \qquad + \int_{\R^d} \mu(\ud \zeta)\left|\mathcal{F}\left(\Sigma^k_{x,z}(s,\cdot) Y(t-s,\cdot)\right)(\zeta)\right|^2  \\
		& =: Q^k_{2,1} + Q^k_{2,2}.
	\end{align*}
As for $Q^k_{2,1}$\footnote{There is a comment made in \cite[Proposition 4.8]{SoleMillet2021} and cite th 3.2 hu huang nualart that minkowsi inequality is being used for the first inequality but it seems incorrect how they apply minkowski's and also unnecessary.},
	\begin{align*}
		\Norm{Q^k_{2,1}}_{p/2} & \le C |\xi|^{2r} \Norm{\sup_{(s,y) \in [0,T]\times \R^d} \Sigma^k_{x,z}(s,y)  \int_0^t \ud s \int_{\R^d} \mu(\ud \zeta) \: |\zeta|^{2r} \left| \mathcal{F}\left(Y(t-s,\cdot)\right)(\zeta)\right|^2}_{p/2} \\
		&\le C |\xi|^{2r} \sup_{(s,y) \in [0,T]\times \R^d} \Norm{\Sigma^k_{x,z}(s,y)}^2_{p} \int_0^t \ud s \int_{\R^d} \mu(\ud \zeta) \: |\zeta|^{2r} \left| \mathcal{F}\left(Y(t-s,\cdot)\right)(\zeta)\right|^2 \\
		&\le C |\xi|^{2r} \left[c(\sigma) + L(\sigma) \sup_{(s,y) \in [0,T]\times \R^d} \Norm{u(s,y)}_{p}\right]^2 \int_0^t \ud s \int_{\R^d} \mu(\ud \zeta) \: |\zeta|^{2r} \left| \mathcal{F}\left(Y(t-s,\cdot)\right)(\zeta)\right|^2 
	\end{align*}
which follows from the global Lipschitz assumption on $\sigma$. Now as for $Q^k_{2,2}$,
	\begin{align*}
		\Norm{Q^k_{2,2}}_{p/2} & \le \int_0^t \ud s \sup_{y \in \R^d}\Norm{\Sigma^k_{x,z}(s,y)}_p^2 \int_{\R^d} \mu(\ud \zeta)\left|\mathcal{F}\left( Y(t-s,\cdot)\right)(\zeta)\right|^2 \\
		& \le C \int_0^t \ud s \sup_{y \in \R^d}\Norm{\Sigma^k_{x,z}(s,y)}_p^2   (t-s)^{2(\beta + \gamma -1)-\beta \lambda / \alpha} \int_{\R^{d}}  \frac{ \mu(\ud \zeta) }{\left(1+ |\zeta|^\alpha\right)^2 } \\
		& \le C \left(\int_{\R^{d}}  \frac{ \mu(\ud \zeta) }{\left(1+ |\zeta|^\alpha\right)^2 } \right)L^2(\sigma) \int_0^t \ud s \:  (t-s)^{2(\beta + \gamma -1)-\beta \lambda / \alpha} \sup_{|z_1-z_2| = |\xi| }\Norm{u(s,z_1) - u(s, z_2)}_p^2  \\
		&\le C \left(\int_{\R^{d}}  \frac{ \mu(\ud \zeta) }{\left(1+ |\zeta|^\alpha\right)^2 } \right)L^2(\sigma) t^{2(\beta + \gamma -1)-\beta \lambda / \alpha}\int_0^t \ud s \:  \sup_{|z_1-z_2| = |\xi| }\Norm{u(s,z_1) - u(s, z_2)}_p^2,
	\end{align*}
	where in the second to last line we applied Lemma \ref{L:YfYbdd} and in the last line we assume that $2(\beta + \gamma -1)-\beta \lambda / \alpha \ge 0$. Summarizing the above gives us that 
	\begin{align*}
		\Norm{Q^k_2}_{p/2} &\le C |\xi|^{2r} \left[c(\sigma) + L(\sigma) \sup_{(s,y) \in [0,T]\times \R^d} \Norm{u(s,y)}_{p}\right]^2 \int_0^t \ud s \int_{\R^d} \mu(\ud \zeta) \: |\zeta|^{2r} \left| \mathcal{F}\left(Y(t-s,\cdot)\right)(\zeta)\right|^2 \\
		& \qquad + C \left(\int_{\R^{d}}  \frac{ \mu(\ud \zeta) }{\left(1+ |\zeta|^\alpha\right)^2 } \right)L^2(\sigma) t^{2(\beta + \gamma -1)-\beta \lambda / \alpha}\int_0^t \ud s \:  \sup_{|z_1-z_2| = |\xi| }\Norm{u(s,z_1) - u(s, z_2)}_p^2.
	\end{align*}
	Also by swapping $x$ and $z$ and noticing that $|e^{-i\xi \zeta} - 1| = |e^{i \xi \zeta} - 1|$ then the bound for $ \Norm{Q^k_3}_{p/2} $ with be the same as for $ \Norm{Q^k_2}_{p/2} $.
	
	Lastly, for $Q^k_4$ we have that after applying Lemma \ref{L:Y_phiParseval} with $\phi = \psi = \sum^k_x(s,\cdot)$,
	\begin{align*}
		|Q^k_4| &\le C \int_0^t \ud s \int_{\R^d} \mu(\ud \zeta) \left|2 - e^{-i\xi \zeta} - e^{i\xi \zeta} \right| \left| \mathcal{F}\left(\Sigma^k_x(s,\cdot) Y(t-s,\cdot)\right)(\zeta) \right|^2 \\
		&\le C |\xi|^{2r} \int_0^t \ud s \int_{\R^d} \mu(\ud \zeta) |\zeta|^{2r} \left| \mathcal{F}\left(\Sigma^k_x(s,\cdot) Y(t-s,\cdot)\right)(\zeta) \right|^2 \\
		&= C |\xi|^{2r} \int_0^t \ud s \int_{\R^d} \ud y \: g(y) \left[\left(\Sigma^k_x(s,\cdot) Y(t-s,\cdot)\right)*\left(\widetilde{\Sigma^k_x(s,\cdot) Y(t-s,\cdot)}\right)\right](y),
	\end{align*}
where $g = \mathcal{F}^{-1} \left(\mu(\cdot) |\cdot|^{2r} \right)$. Thus we see that 
	\begin{align*}
			\Norm{Q^k_4}_{p/2} &\le C |\xi|^{2r} \int_0^t \ud s \: \sup_{y \in \R^d} \Norm{\Sigma^k_x(s,y)}_{p}^2 \int_{\R^d} \ud y \: g(y) \left[\left( Y(t-s,\cdot)\right)*\left(\widetilde{ Y(t-s,\cdot)}\right)\right](y) \\
			&= C|\xi|^{2r} \sup_{(s,y) \in [0,T] \times \R^d} \Norm{\Sigma^k_x(s,y)}_{p}^2 \int_0^t \ud s \: \int_{\R^d} \mu(\ud \zeta) |\zeta|^{2r} \left| \mathcal{F}\left( Y(t-s,\cdot)\right)(\zeta)\right|^2 \\
			&\le C|\xi|^{2r} \left[c(\sigma) + L(\sigma)\sup_{(s,y) \in [0,T] \times \R^d} \Norm{u(s,y)}_{p}\right]^2 \int_0^t \ud s \: \int_{\R^d} \mu(\ud \zeta) |\zeta|^{2r} \left| \mathcal{F}\left( Y(t-s,\cdot)\right)(\zeta)\right|^2.
	\end{align*}
Thus taking \eqref{E:Ispaceinc} into consideration, we have that 
	\begin{align*}
		\Norm{I(t,x) - I(t,z)}_p^2 &\le C|\xi|^{2r} \left[c(\sigma) + L(\sigma)\sup_{(s,y) \in [0,T] \times \R^d} \Norm{u(s,y)}_{p}\right]^2 \int_0^t \ud s \: \int_{\R^d} \mu(\ud \zeta) |\zeta|^{2r} \left| \mathcal{F}\left( Y(t-s,\cdot)\right)(\zeta)\right|^2 \\
		& \qquad  + C \left(\int_{\R^{d}}  \frac{ \mu(\ud \zeta) }{\left(1+ |\zeta|^\alpha\right)^2 } \right)L^2(\sigma) t^{2(\beta + \gamma -1)-\beta \lambda / \alpha}\int_0^t \ud s \:  \sup_{|z_1-z_2| = |\xi| }\Norm{u(s,z_1) - u(s, z_2)}_p^2. 
	\end{align*}
Lastly we note that by the same techniques in Lemma \ref{L:YfYbdd}, we see that
	\begin{align*}
		\int_0^t \ud s \: \int_{\R^d} \mu(\ud \zeta) |\zeta|^{2r} & \left| \mathcal{F}\left( Y(t-s,\cdot)\right)(\zeta)\right|^2 \\
		 &\le \int_0^t \ud s \: (t-s)^{2(\beta + \gamma -1)}\int_{\R^d} \mu(\ud \zeta) |\zeta|^{2r} E_{\beta, \beta + \gamma}\left(-2^{-1}\nu(t-s)^{\beta} |\zeta|^{\alpha}\right)^2 \\
		 &\le C \int_0^t \ud s \: (t-s)^{2(\beta + \gamma -1)}\int_{\R^d} \mu(\ud \zeta) |\zeta|^{2r} \left(1+2^{-1}\nu(t-s)^{\beta} |\zeta|^{\alpha}\right)^{-2} \\
		 &= C \int_0^t \ud s \: (t-s)^{2(\beta + \gamma -1) - \frac{\beta \lambda}{\alpha} - \frac{2r\beta}{\alpha}}\int_{\R^d} \mu(\ud \zeta) \frac{|\zeta|^{2r}}{\left(1+ |\zeta|^{\alpha}\right)^{2}} ,
	\end{align*}
therefore,
	\begin{align*}
&\Norm{I(t,x) - I(t,z)}_p^2 \\
&\le C|\xi|^{2r} \left[\int_{\R^d} \mu(\ud \zeta) \frac{|\zeta|^{2r}}{\left(1+ |\zeta|^{\alpha}\right)^{2}} \right] \left[\int_0^t \ud s \:   (t-s)^{2(\beta + \gamma -1) - \frac{\beta \lambda}{\alpha} - \frac{2r\beta}{\alpha}}\right] \left[c(\sigma) + L(\sigma)\sup_{(s,y) \in [0,T] \times \R^d} \Norm{u(s,y)}_{p}\right]^2 \\
& \qquad  + C \left[\int_{\R^{d}}  \frac{ \mu(\ud \zeta) }{\left(1+ |\zeta|^\alpha\right)^2 } \right]L^2(\sigma) t^{2(\beta + \gamma -1)-\beta \lambda / \alpha}\int_0^t \ud s \:  \sup_{|z_1-z_2| = |\xi| }\Norm{u(s,z_1) - u(s, z_2)}_p^2,
\end{align*}
which requires us to assume that 
	\[
	2(\beta + \gamma -1) - \frac{\beta \lambda}{\alpha} >0
	\quad \text{and} \quad
	2(\beta + \gamma -1) - \frac{\beta \lambda}{\alpha} - \frac{2r\beta}{\alpha} > -1
	\]

\end{proof}

	\addcontentsline{toc}{section}{Bibliography}
	\begin{thebibliography}{999}
		%sample
		%\bibitem{Chen17Tran}% {{{ 3.
		%Chen, L. (2017).
		%\newblock Nonlinear stochastic time-fractional diffusion equations on $\R$: moments, H\"older regularity and intermittency.
		%\newblock {\em Trans. Amer. Math. Soc.} \textbf{369} (12) 8497--8535.
		\bibitem{CH2017}
		Chen L. and Huang J. (2017).
		\newblock Comparison Principle for SHE on $\R^d$.
		\newblock {\em The Annals of Probability} \textbf{Vol. 47, No. 2} 989 -- 1035.
		
		\bibitem{ChenHu21}
		Chen, L and Hu, G. (2021).
		\newblock H\"older regularity of the nonlinear stochastic time-fractional slow and fast diffusion equations on $\R^d$.
		\newblock{\em arXiv:2105.00891.}
		
		\bibitem{CHHH16}
		Chen, L., Hu, G., Hu, Y. and Huang, J. (2016)
		\newblock{Space-time fractional diffusions in Gaussian noisy environment}
		\newblock{\it Stochastics}, 89:1, 171--206.
		
		\bibitem{CHN19}% {{{ 5.
		Chen, L., Hu, Y., and Nualart, D. (2019).
		\newblock Nonlinear stochastic time-fractional slow and fast diffusion equations on $\R^d$.
		\newblock {\em Stochastic Process. Appl.} \textbf{129}, 5073--5112
		% }}}
		
		\bibitem{FN2016}
		Foondun, M and Nane E. (2016)
		\newblock{Asymptotic properties of some space-time fractional stochastic equations.}
		\newblock{\em Mathematische Zeitschrift} 2015, \textbf{287}, 493 -- 519
		
		\bibitem{FT2014}
		Foondun, M. and Tian, K. (2014)
		\newblock On some properties of a class of fractional stochastic heat equations.
		\newblock {\em arXiv:1404.6791v1}
		
		\bibitem{HHN13}
		Hu, Y., Huang, J. and Nualart, D. (2013)
		\newblock On H\"older continuity of the solution of stochastic wave
		equations
		\newblock {\it arXiv:1308.6607v1}
		
		\bibitem{Kilbas}
		\newblock H-Transforms: Theory and Applications 
		\newblock{\it An International Series of Monographs in Mathematics.}
		
		\bibitem{Kuo}
		Kuo, H. (2000)
		\newblock{\em Introduction to Stochastic Integration}
		\newblock Springer, Universitext
		
		\bibitem{Podlubny99}
		Podlubny, I. (1999).
		\newblock{\em Fractional Differential Equations.}
		\newblock Mathematics in Science and Engineering, Academic Press, Vol198 (1999).
		
		\bibitem{SoleMillet2021}
		Sanz-Sole, M. and Millet, A. (2021)
		\newblock{Global solutions to stochastic wave equations with superlinear coefficients.}
		\newblock {\em Stochastic Processes and their Applications}, Elsevier, 2021, 139, pp.175-211.
		
		\bibitem{Walsh89}
		Walsh, J. (1989)
		\newblock{\em An Introduction to Stochastic Partial Differential Equations.}
		\newblock \'Ecole d'\'et\'e de Probabilit\'es de Saint-Flour, XIV|1994. Lecture Notes in Math., vol. 1180, pp. 265-439. Springer, Berlin (1986).
	\end{thebibliography}
	
\end{document}
